% =============================================================================
% pt_tikz_styles.tex --- Process-tensor graphical calculus (SciPost / docs)
% =============================================================================
%
% Visual conventions
% ------------------
% Time in every process-tensor diagram runs RIGHT TO LEFT:
%
%     end  <----------------  start
%               time
%     t_n                     t_0
%
% Earlier operations sit to the right of later ones.  The reusable time arrow
% always points left.  Never reverse this convention in a panel.
%
% Spaces are distinguished by LINE STRUCTURE, not colour (wires are black):
%   * Hilbert ket leg     --- solid wire.
%   * Hilbert bra leg     --- solid wire plus an optional prime marker.
%   * Liouville leg       --- double wire (fused s \otimes s^*, dim d^2).
%   * PT-MPO link index   --- solid wire at twice the site-leg thickness.
%   * Ordinary TN bond    --- solid wire at the site-leg thickness.
%   * Classical outcome   --- dashed wire.
%
% Boundary states are equilateral triangles, not arrowheads:
%   * Ket  ----|>     right-pointing triangle; wire enters from the left.
%   * Bra  <|----     left-pointing triangle;  wire exits to the right.
% Outlines are black.  Fills may differ (ket / bra / Liouville / instrument).
% A Liouville vector has a plain triangle; only its leg is double-wired.
%
% Operators and two-legged tensors are squares with a thick stroke and slight
% corner rounding.  Process cores use a light-blue fill; instruments use a
% light-gold fill; ordinary operators use a light-grey fill.
% Filled shapes are drawn on a foreground layer so wires never show through.
% Trace symbols are cups (Hilbert) or compact Tr nodes (Liouville), never
% ordinary measurement boxes.
%
% Input / output ports on a process core (local time also right-to-left):
%   * Output of the process (to the instrument) sits on the RIGHT of the pair.
%   * Input  of the process (from the instrument) sits on the LEFT of the pair.
%
% Column-major vectorisation (Julia / ProcessTensors.jl):
%     |A>> = sum_{j,k} A_{jk} |k> \otimes |j>
%     vec(|j><k|) = |k> \otimes |j>
%     vec(A \rho B) = (B^T \otimes A) |\rho>>
%
% Public commands (all prefixed PT)
% ---------------------------------
%   \PTKet / \PTBra / \PTLiouvilleKet / \PTLiouvilleBra
%   \PTTensor / \PTOperator / \PTSuperoperator / \PTProcessCore / \PTInstrument
%   \PTInstrumentNode          generic instrument constructor
%   \PTIdentity / \PTControl / \PTUnitary / \PTPrepare / \PTMeasure
%   \PTProjective / \PTCausalBreak / \PTTraceHilbert / \PTTraceLiouville
%   \PTCombiner / \PTCombinerInv
%   \PTIndexLabel              math + ITensor tag + optional prime
%   \PTCup / \PTCap / \PTTimeArrow
%   \PTDrawProcessMPO[n=4, at={0,0}, name=pt]  % at = x,y without extra parens
%   \PTDrawJointProcess        wrapper with joint=true
%
% Introducing a new instrument
% ----------------------------
% Derive it from \PTInstrumentNode (or from a ket/bra/trace kind).  Example:
%
%   \newcommand{\PTReset}[4][]{%
%     \PTInstrumentNode[fill=PTInstrument, classical=none, #1]{#2}{#3}{#4}%
%   }
%   % usage: \PTReset{R}{(0,0)}{\(\mathcal{R}_k\)}
%
% Future instruments (testers, feedback, multi-site, coarse-grained
% measurements, discard-and-reprepare, ...) should likewise be thin wrappers
% around \PTInstrumentNode so existing figures keep compiling.
%
% Overriding colours
% ------------------
% All styles refer to named colours.  Redefine them AFTER inputting this
% file and BEFORE any tikzpicture; do not edit the drawing commands:
%
%   \definecolor{PTLink}{HTML}{111111}
%   \definecolor{PTProcess}{HTML}{F0F0F0}
%
% Lengths (\PTCoreWidth, \PTTimeSpacing, ...) may be \setlength similarly.
%
% Loading
% -------
% Preamble only:
%   % =============================================================================
% pt_tikz_styles.tex --- Process-tensor graphical calculus (SciPost / docs)
% =============================================================================
%
% Visual conventions
% ------------------
% Time in every process-tensor diagram runs RIGHT TO LEFT:
%
%     end  <----------------  start
%               time
%     t_n                     t_0
%
% Earlier operations sit to the right of later ones.  The reusable time arrow
% always points left.  Never reverse this convention in a panel.
%
% Spaces are distinguished by LINE STRUCTURE, not colour (wires are black):
%   * Hilbert ket leg     --- solid wire.
%   * Hilbert bra leg     --- solid wire plus an optional prime marker.
%   * Liouville leg       --- double wire (fused s \otimes s^*, dim d^2).
%   * PT-MPO link index   --- solid wire at twice the site-leg thickness.
%   * Ordinary TN bond    --- solid wire at the site-leg thickness.
%   * Classical outcome   --- dashed wire.
%
% Boundary states are equilateral triangles, not arrowheads:
%   * Ket  ----|>     right-pointing triangle; wire enters from the left.
%   * Bra  <|----     left-pointing triangle;  wire exits to the right.
% Outlines are black.  Fills may differ (ket / bra / Liouville / instrument).
% A Liouville vector has a plain triangle; only its leg is double-wired.
%
% Operators and two-legged tensors are squares with a thick stroke and slight
% corner rounding.  Process cores use a light-blue fill; instruments use a
% light-gold fill; ordinary operators use a light-grey fill.
% Filled shapes are drawn on a foreground layer so wires never show through.
% Trace symbols are cups (Hilbert) or compact Tr nodes (Liouville), never
% ordinary measurement boxes.
%
% Input / output ports on a process core (local time also right-to-left):
%   * Output of the process (to the instrument) sits on the RIGHT of the pair.
%   * Input  of the process (from the instrument) sits on the LEFT of the pair.
%
% Column-major vectorisation (Julia / ProcessTensors.jl):
%     |A>> = sum_{j,k} A_{jk} |k> \otimes |j>
%     vec(|j><k|) = |k> \otimes |j>
%     vec(A \rho B) = (B^T \otimes A) |\rho>>
%
% Public commands (all prefixed PT)
% ---------------------------------
%   \PTKet / \PTBra / \PTLiouvilleKet / \PTLiouvilleBra
%   \PTTensor / \PTOperator / \PTSuperoperator / \PTProcessCore / \PTInstrument
%   \PTInstrumentNode          generic instrument constructor
%   \PTIdentity / \PTControl / \PTUnitary / \PTPrepare / \PTMeasure
%   \PTProjective / \PTCausalBreak / \PTTraceHilbert / \PTTraceLiouville
%   \PTCombiner / \PTCombinerInv
%   \PTIndexLabel              math + ITensor tag + optional prime
%   \PTCup / \PTCap / \PTTimeArrow
%   \PTDrawProcessMPO[n=4, at={0,0}, name=pt]  % at = x,y without extra parens
%   \PTDrawJointProcess        wrapper with joint=true
%
% Introducing a new instrument
% ----------------------------
% Derive it from \PTInstrumentNode (or from a ket/bra/trace kind).  Example:
%
%   \newcommand{\PTReset}[4][]{%
%     \PTInstrumentNode[fill=PTInstrument, classical=none, #1]{#2}{#3}{#4}%
%   }
%   % usage: \PTReset{R}{(0,0)}{\(\mathcal{R}_k\)}
%
% Future instruments (testers, feedback, multi-site, coarse-grained
% measurements, discard-and-reprepare, ...) should likewise be thin wrappers
% around \PTInstrumentNode so existing figures keep compiling.
%
% Overriding colours
% ------------------
% All styles refer to named colours.  Redefine them AFTER inputting this
% file and BEFORE any tikzpicture; do not edit the drawing commands:
%
%   \definecolor{PTLink}{HTML}{111111}
%   \definecolor{PTProcess}{HTML}{F0F0F0}
%
% Lengths (\PTCoreWidth, \PTTimeSpacing, ...) may be \setlength similarly.
%
% Loading
% -------
% Preamble only:
%   % =============================================================================
% pt_tikz_styles.tex --- Process-tensor graphical calculus (SciPost / docs)
% =============================================================================
%
% Visual conventions
% ------------------
% Time in every process-tensor diagram runs RIGHT TO LEFT:
%
%     end  <----------------  start
%               time
%     t_n                     t_0
%
% Earlier operations sit to the right of later ones.  The reusable time arrow
% always points left.  Never reverse this convention in a panel.
%
% Spaces are distinguished by LINE STRUCTURE, not colour (wires are black):
%   * Hilbert ket leg     --- solid wire.
%   * Hilbert bra leg     --- solid wire plus an optional prime marker.
%   * Liouville leg       --- double wire (fused s \otimes s^*, dim d^2).
%   * PT-MPO link index   --- solid wire at twice the site-leg thickness.
%   * Ordinary TN bond    --- solid wire at the site-leg thickness.
%   * Classical outcome   --- dashed wire.
%
% Boundary states are equilateral triangles, not arrowheads:
%   * Ket  ----|>     right-pointing triangle; wire enters from the left.
%   * Bra  <|----     left-pointing triangle;  wire exits to the right.
% Outlines are black.  Fills may differ (ket / bra / Liouville / instrument).
% A Liouville vector has a plain triangle; only its leg is double-wired.
%
% Operators and two-legged tensors are squares with a thick stroke and slight
% corner rounding.  Process cores use a light-blue fill; instruments use a
% light-gold fill; ordinary operators use a light-grey fill.
% Filled shapes are drawn on a foreground layer so wires never show through.
% Trace symbols are cups (Hilbert) or compact Tr nodes (Liouville), never
% ordinary measurement boxes.
%
% Input / output ports on a process core (local time also right-to-left):
%   * Output of the process (to the instrument) sits on the RIGHT of the pair.
%   * Input  of the process (from the instrument) sits on the LEFT of the pair.
%
% Column-major vectorisation (Julia / ProcessTensors.jl):
%     |A>> = sum_{j,k} A_{jk} |k> \otimes |j>
%     vec(|j><k|) = |k> \otimes |j>
%     vec(A \rho B) = (B^T \otimes A) |\rho>>
%
% Public commands (all prefixed PT)
% ---------------------------------
%   \PTKet / \PTBra / \PTLiouvilleKet / \PTLiouvilleBra
%   \PTTensor / \PTOperator / \PTSuperoperator / \PTProcessCore / \PTInstrument
%   \PTInstrumentNode          generic instrument constructor
%   \PTIdentity / \PTControl / \PTUnitary / \PTPrepare / \PTMeasure
%   \PTProjective / \PTCausalBreak / \PTTraceHilbert / \PTTraceLiouville
%   \PTCombiner / \PTCombinerInv
%   \PTIndexLabel              math + ITensor tag + optional prime
%   \PTCup / \PTCap / \PTTimeArrow
%   \PTDrawProcessMPO[n=4, at={0,0}, name=pt]  % at = x,y without extra parens
%   \PTDrawJointProcess        wrapper with joint=true
%
% Introducing a new instrument
% ----------------------------
% Derive it from \PTInstrumentNode (or from a ket/bra/trace kind).  Example:
%
%   \newcommand{\PTReset}[4][]{%
%     \PTInstrumentNode[fill=PTInstrument, classical=none, #1]{#2}{#3}{#4}%
%   }
%   % usage: \PTReset{R}{(0,0)}{\(\mathcal{R}_k\)}
%
% Future instruments (testers, feedback, multi-site, coarse-grained
% measurements, discard-and-reprepare, ...) should likewise be thin wrappers
% around \PTInstrumentNode so existing figures keep compiling.
%
% Overriding colours
% ------------------
% All styles refer to named colours.  Redefine them AFTER inputting this
% file and BEFORE any tikzpicture; do not edit the drawing commands:
%
%   \definecolor{PTLink}{HTML}{111111}
%   \definecolor{PTProcess}{HTML}{F0F0F0}
%
% Lengths (\PTCoreWidth, \PTTimeSpacing, ...) may be \setlength similarly.
%
% Loading
% -------
% Preamble only:
%   % =============================================================================
% pt_tikz_styles.tex --- Process-tensor graphical calculus (SciPost / docs)
% =============================================================================
%
% Visual conventions
% ------------------
% Time in every process-tensor diagram runs RIGHT TO LEFT:
%
%     end  <----------------  start
%               time
%     t_n                     t_0
%
% Earlier operations sit to the right of later ones.  The reusable time arrow
% always points left.  Never reverse this convention in a panel.
%
% Spaces are distinguished by LINE STRUCTURE, not colour (wires are black):
%   * Hilbert ket leg     --- solid wire.
%   * Hilbert bra leg     --- solid wire plus an optional prime marker.
%   * Liouville leg       --- double wire (fused s \otimes s^*, dim d^2).
%   * PT-MPO link index   --- solid wire at twice the site-leg thickness.
%   * Ordinary TN bond    --- solid wire at the site-leg thickness.
%   * Classical outcome   --- dashed wire.
%
% Boundary states are equilateral triangles, not arrowheads:
%   * Ket  ----|>     right-pointing triangle; wire enters from the left.
%   * Bra  <|----     left-pointing triangle;  wire exits to the right.
% Outlines are black.  Fills may differ (ket / bra / Liouville / instrument).
% A Liouville vector has a plain triangle; only its leg is double-wired.
%
% Operators and two-legged tensors are squares with a thick stroke and slight
% corner rounding.  Process cores use a light-blue fill; instruments use a
% light-gold fill; ordinary operators use a light-grey fill.
% Filled shapes are drawn on a foreground layer so wires never show through.
% Trace symbols are cups (Hilbert) or compact Tr nodes (Liouville), never
% ordinary measurement boxes.
%
% Input / output ports on a process core (local time also right-to-left):
%   * Output of the process (to the instrument) sits on the RIGHT of the pair.
%   * Input  of the process (from the instrument) sits on the LEFT of the pair.
%
% Column-major vectorisation (Julia / ProcessTensors.jl):
%     |A>> = sum_{j,k} A_{jk} |k> \otimes |j>
%     vec(|j><k|) = |k> \otimes |j>
%     vec(A \rho B) = (B^T \otimes A) |\rho>>
%
% Public commands (all prefixed PT)
% ---------------------------------
%   \PTKet / \PTBra / \PTLiouvilleKet / \PTLiouvilleBra
%   \PTTensor / \PTOperator / \PTSuperoperator / \PTProcessCore / \PTInstrument
%   \PTInstrumentNode          generic instrument constructor
%   \PTIdentity / \PTControl / \PTUnitary / \PTPrepare / \PTMeasure
%   \PTProjective / \PTCausalBreak / \PTTraceHilbert / \PTTraceLiouville
%   \PTCombiner / \PTCombinerInv
%   \PTIndexLabel              math + ITensor tag + optional prime
%   \PTCup / \PTCap / \PTTimeArrow
%   \PTDrawProcessMPO[n=4, at={0,0}, name=pt]  % at = x,y without extra parens
%   \PTDrawJointProcess        wrapper with joint=true
%
% Introducing a new instrument
% ----------------------------
% Derive it from \PTInstrumentNode (or from a ket/bra/trace kind).  Example:
%
%   \newcommand{\PTReset}[4][]{%
%     \PTInstrumentNode[fill=PTInstrument, classical=none, #1]{#2}{#3}{#4}%
%   }
%   % usage: \PTReset{R}{(0,0)}{\(\mathcal{R}_k\)}
%
% Future instruments (testers, feedback, multi-site, coarse-grained
% measurements, discard-and-reprepare, ...) should likewise be thin wrappers
% around \PTInstrumentNode so existing figures keep compiling.
%
% Overriding colours
% ------------------
% All styles refer to named colours.  Redefine them AFTER inputting this
% file and BEFORE any tikzpicture; do not edit the drawing commands:
%
%   \definecolor{PTLink}{HTML}{111111}
%   \definecolor{PTProcess}{HTML}{F0F0F0}
%
% Lengths (\PTCoreWidth, \PTTimeSpacing, ...) may be \setlength similarly.
%
% Loading
% -------
% Preamble only:
%   \input{figures/tikz/pt_tikz_styles.tex}
% Figure fragments then:
%   \input{figures/tikz/<name>.tikz}
%
% =============================================================================

\ifdefined\PTTikZLoaded
  \expandafter\endinput
\fi
\def\PTTikZLoaded{1}

\RequirePackage{tikz}
\RequirePackage{xcolor}
\RequirePackage{keyval}
\usetikzlibrary{
  arrows.meta,
  backgrounds,
  calc,
  decorations.markings,
  fit,
  matrix,
  positioning,
  shapes.geometric
}

% -----------------------------------------------------------------------------
% Colours (Okabe--Ito / colour-blind friendly; intelligible in grayscale)
% -----------------------------------------------------------------------------
\definecolor{PTHilbertKet}{HTML}{0072B2}
\definecolor{PTHilbertBra}{HTML}{D55E00}
\definecolor{PTLiouville}{HTML}{7B3294}
\definecolor{PTLink}{HTML}{009E73}
\definecolor{PTTensor}{HTML}{E8E8E8}
\definecolor{PTProcess}{HTML}{DCEAF7}
\definecolor{PTInstrument}{HTML}{F6E8A6}
\definecolor{PTTrace}{HTML}{CC79A7}
\definecolor{PTMeasurement}{HTML}{E69F00}
\definecolor{PTLabel}{HTML}{333333}
\definecolor{PTBond}{HTML}{6A6A6A}

% -----------------------------------------------------------------------------
% Lengths (the single place to retune the visual language)
% -----------------------------------------------------------------------------
\newlength{\PTCoreWidth}      \setlength{\PTCoreWidth}{9.6mm}
\newlength{\PTCoreHeight}     \setlength{\PTCoreHeight}{9.6mm}
\newlength{\PTTimeSpacing}    \setlength{\PTTimeSpacing}{20.0mm}
\newlength{\PTLegLength}      \setlength{\PTLegLength}{7.0mm}
\newlength{\PTIOSep}          \setlength{\PTIOSep}{2.4mm}
% Equilateral triangle: vertical base = side, horizontal altitude = (sqrt(3)/2)*side.
\newlength{\PTKetHeight}      \setlength{\PTKetHeight}{7.8mm}
\newlength{\PTKetWidth}       \setlength{\PTKetWidth}{6.75mm}
\newlength{\PTBendDepth}      \setlength{\PTBendDepth}{5.5mm}
\newlength{\PTCombinerSize}   \setlength{\PTCombinerSize}{7.2mm}
\newlength{\PTHilbertLW}      \setlength{\PTHilbertLW}{1.35pt}
\newlength{\PTLiouvilleLW}    \setlength{\PTLiouvilleLW}{0.70pt}
\newlength{\PTLiouvilleGap}   \setlength{\PTLiouvilleGap}{1.55pt}
\newlength{\PTLinkLW}         \setlength{\PTLinkLW}{2.70pt}
\newlength{\PTBondLW}         \setlength{\PTBondLW}{1.35pt}
\newlength{\PTClassicalLW}    \setlength{\PTClassicalLW}{1.15pt}
\newlength{\PTCoreStroke}     \setlength{\PTCoreStroke}{1.45pt}
\newlength{\PTCoreRadius}     \setlength{\PTCoreRadius}{1.3pt}
\newlength{\PTBridgeHalo}     \setlength{\PTBridgeHalo}{4.2pt}

% Scratch length for the MPO drawer (not a public tuning knob).
\newlength{\PTMPOspacing}

% -----------------------------------------------------------------------------
% Layers: wires behind, labels in the middle, filled shapes in front.
% -----------------------------------------------------------------------------
\pgfdeclarelayer{PTforeground}
\pgfsetlayers{background,main,PTforeground}
\tikzset{
  on PTforeground layer/.style={
    execute at begin scope={\pgfonlayer{PTforeground}},
    execute at end scope={\endpgfonlayer},
  },
}

% -----------------------------------------------------------------------------
% Line and node styles
% -----------------------------------------------------------------------------
\tikzset{
  PTPicture/.style={
    baseline,
    line join=round,
    >={Stealth[length=2.0mm, width=1.45mm]},
  },
  % --- wires (black; structure, not colour, distinguishes spaces) ----------
  PTWire/.style={
    draw=black,
    line width=\PTHilbertLW,
    line cap=butt,
    line join=round,
    double=none,
  },
  PTHilbertKetLeg/.style={
    PTWire,
    on background layer,
  },
  PTHilbertBraLeg/.style={
    draw=black,
    line width=\PTHilbertLW,
    line cap=butt,
    line join=round,
    double=none,
    on background layer,
  },
  % Optional prime tick at the midpoint of a bra wire.
  PTHilbertBraLeg primed/.style={
    PTHilbertBraLeg,
    postaction={
      decorate,
      decoration={
        markings,
        mark=at position 0.55 with {
          \node[font=\scriptsize, text=PTLabel, inner sep=0.4pt,
                xshift=3.2pt, yshift=0.4pt] {\(\prime\)};
        }
      }
    },
  },
  PTLiouvilleLeg/.style={
    draw=black,
    line width=\PTLiouvilleLW,
    double=white,
    double distance=\PTLiouvilleGap,
    line cap=butt,
    line join=round,
    on background layer,
  },
  PTBondLeg/.style={
    draw=black,
    line width=\PTBondLW,
    line cap=butt,
    line join=round,
    double=none,
    on background layer,
  },
  PTLinkLeg/.style={
    draw=black,
    line width=\PTLinkLW,
    line cap=butt,
    line join=round,
    double=none,
    on background layer,
  },
  PTClassicalLeg/.style={
    draw=black,
    line width=\PTClassicalLW,
    dash pattern=on 2.4pt off 1.6pt,
    line cap=butt,
    line join=round,
    double=none,
    on background layer,
  },
  PTTraceCup/.style={
    draw=PTTrace,
    line width=\PTHilbertLW,
    line cap=round,
    line join=round,
    double=none,
    on background layer,
  },
  % White halo so an over-crossing does not look like a contraction.
  PTBridge/.style={
    preaction={
      draw=white,
      line width=\PTBridgeHalo,
      double=none,
      line cap=butt,
    },
  },
  % --- cores ---------------------------------------------------------------
  PTCore/.style={
    draw=black,
    fill=PTTensor,
    rounded corners=\PTCoreRadius,
    minimum width=\PTCoreWidth,
    minimum height=\PTCoreHeight,
    inner sep=2.2pt,
    outer sep=0pt,
    line width=\PTCoreStroke,
    font=\small,
    text=PTLabel,
    align=center,
  },
  PTTraceNode/.style={
    draw=black,
    fill=PTTrace!18,
    rounded corners=\PTCoreRadius,
    minimum width=7.2mm,
    minimum height=7.2mm,
    inner sep=1.4pt,
    outer sep=0pt,
    line width=\PTCoreStroke,
    font=\small,
    text=PTLabel,
    align=center,
  },
  PTDummyBond/.style={
    circle,
    fill=black,
    draw=none,
    inner sep=1.15pt,
    outer sep=0pt,
  },
}

% -----------------------------------------------------------------------------
% Index labels: math + ITensor tag + optional prime level
% -----------------------------------------------------------------------------
\makeatletter
\define@key{ptindex}{anchor}{\def\PTIndexAnchor{#1}}
\define@key{ptindex}{math}{\def\PTIndexMath{#1}}
\define@key{ptindex}{tag}{\def\PTIndexTag{#1}}
\define@key{ptindex}{prime}{\def\PTIndexPrime{#1}}
\define@key{ptindex}{xshift}{\def\PTIndexX{#1}}
\define@key{ptindex}{yshift}{\def\PTIndexY{#1}}
\makeatother

% Append a visual prime level to math (') and to a tag (ASCII apostrophe).
\newcommand{\PTIndexPrimeTokens}{%
  \ifx\PTIndexPrime\empty
  \else
    \ifnum\PTIndexPrime>0 '\fi
    \ifnum\PTIndexPrime>1 '\fi
    \ifnum\PTIndexPrime>2 '\fi
  \fi
}

\newcommand{\PTIndexLabel}[2][]{%
  \def\PTIndexAnchor{south}%
  \def\PTIndexMath{}%
  \def\PTIndexTag{}%
  \def\PTIndexPrime{}%
  \def\PTIndexX{0pt}%
  \def\PTIndexY{0pt}%
  \setkeys{ptindex}{#1}%
  \node[
    inner sep=0.9pt,
    outer sep=1.2pt,
    anchor=\PTIndexAnchor,
    align=center,
    xshift=\PTIndexX,
    yshift=\PTIndexY,
    text=PTLabel,
  ] at (#2) {%
    \def\PTIndexMathEmpty{}%
    \ifx\PTIndexMath\PTIndexMathEmpty
    \else
      {\small\(\PTIndexMath\PTIndexPrimeTokens\)}%
      \ifx\PTIndexTag\PTIndexMathEmpty
      \else \\[-0.15ex]%
      \fi
    \fi
    \ifx\PTIndexTag\PTIndexMathEmpty
    \else
      {\scriptsize\ttfamily \PTIndexTag\PTIndexPrimeTokens}%
    \fi
  };%
}

% -----------------------------------------------------------------------------
% Ket and bra boundary tensors
% -----------------------------------------------------------------------------
% Invisible rectangular node provides named anchors.  The triangle is drawn
% on top so that .west (ket) / .east (bra) is the wire attachment (base).

\newcommand{\PTKet}[3]{%
  \node[
    name=#1,
    inner sep=0pt,
    outer sep=0pt,
    minimum width=\PTKetWidth,
    minimum height=\PTKetHeight,
    font=\small,
  ] at #2 {};
  \begin{scope}[on PTforeground layer]
    \filldraw[
      draw=black,
      fill=PTHilbertKet!18,
      line width=\PTCoreStroke,
      line join=round,
    ] (#1.west |- #1.south) -- (#1.east) -- (#1.west |- #1.north) -- cycle;
    \node[font=\small, text=PTLabel, inner sep=0.4pt]
      at ($ (#1.west)!0.36!(#1.east) $) {#3};
  \end{scope}
}

\newcommand{\PTBra}[3]{%
  \node[
    name=#1,
    inner sep=0pt,
    outer sep=0pt,
    minimum width=\PTKetWidth,
    minimum height=\PTKetHeight,
    font=\small,
  ] at #2 {};
  \begin{scope}[on PTforeground layer]
    \filldraw[
      draw=black,
      fill=PTHilbertBra!18,
      line width=\PTCoreStroke,
      line join=round,
    ] (#1.east |- #1.south) -- (#1.west) -- (#1.east |- #1.north) -- cycle;
    \node[font=\small, text=PTLabel, inner sep=0.4pt]
      at ($ (#1.east)!0.36!(#1.west) $) {#3};
  \end{scope}
}

\newcommand{\PTLiouvilleKet}[3]{%
  \node[
    name=#1,
    inner sep=0pt,
    outer sep=0pt,
    minimum width=\PTKetWidth,
    minimum height=\PTKetHeight,
    font=\small,
  ] at #2 {};
  \begin{scope}[on PTforeground layer]
    \filldraw[
      draw=black,
      fill=PTLiouville!16,
      line width=\PTCoreStroke,
      line join=round,
    ] (#1.west |- #1.south) -- (#1.east) -- (#1.west |- #1.north) -- cycle;
    \node[font=\small, text=PTLabel, inner sep=0.4pt]
      at ($ (#1.west)!0.36!(#1.east) $) {#3};
  \end{scope}
}

\newcommand{\PTLiouvilleBra}[3]{%
  \node[
    name=#1,
    inner sep=0pt,
    outer sep=0pt,
    minimum width=\PTKetWidth,
    minimum height=\PTKetHeight,
    font=\small,
  ] at #2 {};
  \begin{scope}[on PTforeground layer]
    \filldraw[
      draw=black,
      fill=PTLiouville!16,
      line width=\PTCoreStroke,
      line join=round,
    ] (#1.east |- #1.south) -- (#1.west) -- (#1.east |- #1.north) -- cycle;
    \node[font=\small, text=PTLabel, inner sep=0.4pt]
      at ($ (#1.east)!0.36!(#1.west) $) {#3};
  \end{scope}
}

% -----------------------------------------------------------------------------
% Square tensor cores (thick stroke, slight corner rounding)
% -----------------------------------------------------------------------------
\makeatletter
\define@key{pttensor}{fill}{\def\PTTensorFill{#1}}
\define@key{pttensor}{draw}{\def\PTTensorDrawCol{#1}}
\define@key{pttensor}{minimum width}{\def\PTTensorMinW{#1}}
\define@key{pttensor}{minimum height}{\def\PTTensorMinH{#1}}
\define@key{pttensor}{rounded}{\def\PTTensorRnd{#1}}
\define@key{pttensor}{font}{\def\PTTensorFont{#1}}
\makeatother

\newcommand{\PTTensor}[4][]{%
  \def\PTTensorFill{PTTensor}%
  \def\PTTensorDrawCol{black}%
  \def\PTTensorMinW{\PTCoreWidth}%
  \def\PTTensorMinH{\PTCoreHeight}%
  \def\PTTensorRnd{\PTCoreRadius}%
  \def\PTTensorFont{\small}%
  \setkeys{pttensor}{#1}%
  \begin{scope}[on PTforeground layer]
    \node[
      name=#2,
      PTCore,
      fill=\PTTensorFill,
      draw=\PTTensorDrawCol,
      minimum width=\PTTensorMinW,
      minimum height=\PTTensorMinH,
      rounded corners=\PTTensorRnd,
      font=\PTTensorFont,
    ] at #3 {#4};
  \end{scope}
}

\newcommand{\PTOperator}[4][]{%
  \PTTensor[fill=PTTensor, minimum width=\PTCoreWidth, minimum height=\PTCoreWidth, #1]{#2}{#3}{#4}%
}

\newcommand{\PTSuperoperator}[4][]{%
  \PTTensor[
    fill=PTTensor,
    minimum width=\PTCoreWidth,
    minimum height=\PTCoreWidth,
    #1
  ]{#2}{#3}{#4}%
}

\newcommand{\PTProcessCore}[4][]{%
  \PTTensor[
    fill=PTProcess,
    minimum width=\PTCoreWidth,
    minimum height=\PTCoreWidth,
    #1
  ]{#2}{#3}{#4}%
}

\newcommand{\PTInstrument}[4][]{%
  \PTTensor[
    fill=PTInstrument,
    minimum width=\PTCoreWidth,
    minimum height=\PTCoreWidth,
    #1
  ]{#2}{#3}{#4}%
}

% Vertical Hilbert ket (north) and bra (south) stubs on an operator.
\newcommand{\PTOperatorLegs}[1]{%
  \draw[PTHilbertKetLeg] (#1.north) -- ++(0,\PTLegLength) coordinate (#1-ket);
  \draw[PTHilbertBraLeg] (#1.south) -- ++(0,-\PTLegLength) coordinate (#1-bra);
}

% -----------------------------------------------------------------------------
% Combiners: s \otimes s'  -->  \ell    and the inverse
% -----------------------------------------------------------------------------
\makeatletter
\define@key{ptcomb}{direction}{\def\PTCombDir{#1}}
\makeatother
\def\PTCombDirUp{up}
\def\PTCombDirDown{down}
\def\PTCombDirLeft{left}
\def\PTCombDirRight{right}

\newcommand{\PTCombiner}[4][]{%
  \def\PTCombDir{up}%
  \setkeys{ptcomb}{#1}%
  \ifx\PTCombDir\PTCombDirDown
    \def\PTCombRot{180}%
  \else\ifx\PTCombDir\PTCombDirLeft
    \def\PTCombRot{90}%
  \else\ifx\PTCombDir\PTCombDirRight
    \def\PTCombRot{270}%
  \else
    \def\PTCombRot{0}%
  \fi\fi\fi
  \begin{scope}[on PTforeground layer]
    \node[
      name=#2,
      isosceles triangle,
      isosceles triangle apex angle=60,
      shape border rotate=\PTCombRot,
      inner sep=0pt,
      outer sep=0pt,
      minimum width=\PTCombinerSize,
      draw=black,
      fill=white,
      line width=\PTCoreStroke,
      font=\scriptsize,
      text=PTLabel,
    ] at #3 {};
    \node[font=\scriptsize, text=PTLabel, inner sep=0.8pt, anchor=east]
      at ([xshift=-1.6pt]#2.west) {#4};
  \end{scope}
  \coordinate (#2-fused) at (#2.apex);
  \coordinate (#2-ket)   at (#2.left corner);
  \coordinate (#2-bra)   at (#2.right corner);
}

\newcommand{\PTCombinerInv}[4][]{%
  \PTCombiner[direction=down, #1]{#2}{#3}{#4}%
}

% -----------------------------------------------------------------------------
% Cups, caps, time arrow, callouts
% -----------------------------------------------------------------------------
% Cup (union, opens upward): used for traces and vectorisation bends.
% The two endpoints are the open wire tips; control points sit below them.
\newcommand{\PTCup}[3][]{%
  \draw[#1] (#2)
    .. controls ($ (#2) + (0,-\PTBendDepth) $)
            and ($ (#3) + (0,-\PTBendDepth) $) ..
    (#3);
}

% Trace: left and right semicircles joined by a straight bar BELOW the core
% (never behind it).  #2-loopmid is the bar midpoint for a single index tag.
% Drawn on the main layer so the bar is not covered by the core; the core
% still sits in front at the ports.
\newcommand{\PTTraceLoop}[2][PTWire]{%
  \draw[#1]
    (#2.west)
    arc[start angle=90, delta angle=180, radius=4.0mm]
    -- ([yshift=-8.0mm]#2.east)
    arc[start angle=270, delta angle=180, radius=4.0mm];
  \coordinate (#2-loopmid) at ([yshift=-8.0mm]#2.center);
}

% Transpose bends: left (ket) port over the top towards the right, now a
% primed bra (coordinate #1-i); right (bra) port under the bottom towards
% the left, now an unprimed ket (coordinate #1-j).
\newcommand{\PTTransposeBends}[1]{%
  \draw[PTHilbertBraLeg]
    (#1.west)
    arc[start angle=270, delta angle=-90, radius=3.3mm]
    arc[start angle=180, delta angle=-90, radius=3.3mm]
    coordinate (#1-itop);
  \draw[PTHilbertBraLeg primed]
    (#1-itop) -- ++(11.5mm,0) coordinate (#1-i);
  \draw[PTHilbertKetLeg]
    (#1.east)
    arc[start angle=90, delta angle=-90, radius=3.3mm]
    arc[start angle=0, delta angle=-90, radius=3.3mm]
    -- ++(-11.5mm,0) coordinate (#1-j);
}
\newcommand{\PTCap}[3][]{%
  \draw[#1] (#2)
    .. controls ($ (#2) + (0,\PTBendDepth) $)
            and ($ (#3) + (0,\PTBendDepth) $) ..
    (#3);
}

% Sideways cup used when bending a vertical bra wire into a second output.
% #2 = start (usually the operator's south / bra port).
% #3 = end   (the new upward ket-like port for the bra index).
% The bend goes to the LEFT of both points by default (column-major: |k>
% becomes the left factor).
\newcommand{\PTVecBend}[3][]{%
  \draw[#1] (#2)
    .. controls ($ (#2) + (0,-\PTBendDepth) $)
            and ($ (#3) + (0,-\PTBendDepth) $) ..
    (#3);
}

% Time arrow: ALWAYS left-pointing.  #1 = right coord, #2 = left coord.
% Optional first argument is a TikZ node-anchor for the label (default below).
\newcommand{\PTTimeArrow}[4][below=2.4pt]{%
  \draw[
    -{Stealth[length=2.2mm, width=1.55mm]},
    line width=1.05pt,
    draw=PTLabel,
    line cap=round,
  ] (#2) -- (#3)
    node[midway, #1, font=\small, text=PTLabel, inner sep=0.6pt] {#4};
}

\newcommand{\PTPanelLabel}[2]{%
  \node[
    font=\small\bfseries,
    text=PTLabel,
    inner sep=0.4pt,
    anchor=south west,
  ] at #1 {#2};
}

\newcommand{\PTCallout}[3][]{%
  \node[
    font=\scriptsize,
    text=PTLabel,
    inner sep=0.8pt,
    align=left,
    #1,
  ] at (#2) {#3};
}

\newcommand{\PTMapsTo}[2][]{%
  \node[font=\small, text=PTLabel, inner sep=1pt, #1] at #2 {\(\longmapsto\)};
}

% -----------------------------------------------------------------------------
% Instruments
% -----------------------------------------------------------------------------
\newif\ifPTInstClassical
\PTInstClassicalfalse
\def\PTInstKindBox{box}
\def\PTInstKindKet{ket}
\def\PTInstKindBra{bra}
\def\PTInstKindTrace{trace}

\makeatletter
\define@key{ptinst}{fill}{\def\PTInstFill{#1}}
\define@key{ptinst}{draw}{\def\PTInstDraw{#1}}
\define@key{ptinst}{minimum width}{\def\PTInstMinW{#1}}
\define@key{ptinst}{minimum height}{\def\PTInstMinH{#1}}
\define@key{ptinst}{kind}{\def\PTInstKind{#1}}
\define@key{ptinst}{classical}{\def\PTInstClassDir{#1}}
\define@key{ptinst}{classical label}{\def\PTInstClassLab{#1}}
\define@key{ptinst}{classical length}{\def\PTInstClassLen{#1}}
\makeatother

% Generic constructor.  Quantum legs are NOT auto-drawn: connect to the
% node anchors (north, south, east, west, or the triangle bases).
% Classical legs ARE auto-drawn when classical != none.
\newcommand{\PTInstrumentNode}[4][]{%
  \def\PTInstFill{PTInstrument}%
  \def\PTInstDraw{black}%
  \def\PTInstMinW{\PTCoreWidth}%
  \def\PTInstMinH{\PTCoreWidth}%
  \def\PTInstKind{box}%
  \def\PTInstClassDir{none}%
  \def\PTInstClassLab{}%
  \def\PTInstClassLen{6.2mm}%
  \setkeys{ptinst}{#1}%
  \def\PTInstNone{none}%
  \ifx\PTInstKind\PTInstKindKet
    \node[
      name=#2,
      inner sep=0pt,
      outer sep=0pt,
      minimum width=\PTKetWidth,
      minimum height=\PTKetHeight,
    ] at #3 {};
    \begin{scope}[on PTforeground layer]
      \filldraw[
        draw=black,
        fill=\PTInstFill,
        line width=\PTCoreStroke,
        line join=round,
      ] (#2.west |- #2.south) -- (#2.east) -- (#2.west |- #2.north) -- cycle;
      \node[font=\small, text=PTLabel, inner sep=0.4pt]
        at ($ (#2.west)!0.36!(#2.east) $) {#4};
    \end{scope}
  \else
    \ifx\PTInstKind\PTInstKindBra
      \node[
        name=#2,
        inner sep=0pt,
        outer sep=0pt,
        minimum width=\PTKetWidth,
        minimum height=\PTKetHeight,
      ] at #3 {};
      \begin{scope}[on PTforeground layer]
        \filldraw[
          draw=black,
          fill=\PTInstFill,
          line width=\PTCoreStroke,
          line join=round,
        ] (#2.east |- #2.south) -- (#2.west) -- (#2.east |- #2.north) -- cycle;
        \node[font=\small, text=PTLabel, inner sep=0.4pt]
          at ($ (#2.east)!0.36!(#2.west) $) {#4};
      \end{scope}
    \else
      \ifx\PTInstKind\PTInstKindTrace
        \begin{scope}[on PTforeground layer]
          \node[name=#2, PTTraceNode] at #3 {#4};
        \end{scope}
      \else
        \begin{scope}[on PTforeground layer]
          \node[
            name=#2,
            PTCore,
            fill=\PTInstFill,
            draw=\PTInstDraw,
            minimum width=\PTInstMinW,
            minimum height=\PTInstMinH,
          ] at #3 {#4};
        \end{scope}
      \fi
    \fi
  \fi
  \ifx\PTInstClassDir\PTInstNone
  \else
    \def\PTInstDirDown{down}%
    \def\PTInstDirUp{up}%
    \def\PTInstDirLeft{left}%
    \def\PTInstDirRight{right}%
    \ifx\PTInstClassDir\PTInstDirDown
      \draw[PTClassicalLeg] (#2.south) -- ++(0,-\PTInstClassLen)
        coordinate (#2-classical);
    \else\ifx\PTInstClassDir\PTInstDirUp
      \draw[PTClassicalLeg] (#2.north) -- ++(0,\PTInstClassLen)
        coordinate (#2-classical);
    \else\ifx\PTInstClassDir\PTInstDirLeft
      \draw[PTClassicalLeg] (#2.west) -- ++(-\PTInstClassLen,0)
        coordinate (#2-classical);
    \else
      \draw[PTClassicalLeg] (#2.east) -- ++(\PTInstClassLen,0)
        coordinate (#2-classical);
    \fi\fi\fi
    \def\PTInstClassEmpty{}%
    \ifx\PTInstClassLab\PTInstClassEmpty
    \else
      \node[font=\small, text=PTLabel, inner sep=0.8pt, anchor=north]
        at ([yshift=-1.0pt]#2-classical) {\(\PTInstClassLab\)};
    \fi
  \fi
}

\newcommand{\PTIdentity}[4][]{%
  \PTInstrumentNode[fill=PTTensor, #1]{#2}{#3}{#4}%
}

\newcommand{\PTControl}[4][]{%
  \PTInstrumentNode[#1]{#2}{#3}{#4}%
}

\newcommand{\PTUnitary}[4][]{%
  \PTInstrumentNode[#1]{#2}{#3}{#4}%
}

\newcommand{\PTPrepare}[4][]{%
  \PTInstrumentNode[kind=ket, fill=PTInstrument, #1]{#2}{#3}{#4}%
}

\newcommand{\PTMeasure}[4][]{%
  \PTInstrumentNode[
    classical=down,
    classical label={x},
    #1
  ]{#2}{#3}{#4}%
}

\newcommand{\PTProjective}[4][]{%
  \PTInstrumentNode[
    classical=down,
    classical label={x},
    #1
  ]{#2}{#3}{#4}%
}

% Measurement (earlier, right) then preparation (later, left).
% Quantum wire is broken; the classical outcome is drawn separately.
% Nodes: <name>-E, <name>-P, <name>-x.
\newcommand{\PTCausalBreak}[3][]{%
  \begin{scope}[shift={#3}]
    \PTInstrumentNode[
      classical=none,
      minimum width=8.4mm,
      minimum height=8.4mm,
      #1
    ]{#2-E}{(0,0)}{\(E^{(x)}\)}
    \PTPrepare{#2-P}{(-1.65,0)}{\(\mathcal{P}\)}
    \draw[PTClassicalLeg] (#2-E.south) -- ++(0,-5.8mm) coordinate (#2-x);
    \node[font=\small, text=PTLabel, inner sep=0.6pt, anchor=north]
      at ([yshift=-0.8pt]#2-x) {\(x\)};
  \end{scope}
}

% Hilbert-space trace: magenta cup closing ket and bra tips.
\newcommand{\PTTraceHilbert}[2]{%
  \PTCup[PTTraceCup]{#1}{#2}%
  \node[font=\small, text=PTTrace, inner sep=0.5pt, anchor=north]
    at ($ (#1)!0.5!(#2) + (0,-\PTBendDepth-0.5pt) $) {\(\mathrm{Tr}\)};
}

% Liouville-space trace: compact Tr node sitting on a fused leg.
\newcommand{\PTTraceLiouville}[3]{%
  \begin{scope}[on PTforeground layer]
    \node[name=#1, PTTraceNode] at #2 {#3};
  \end{scope}
}

% -----------------------------------------------------------------------------
% PT-MPO / joint process tensor
% -----------------------------------------------------------------------------
\newif\ifPTMPOjoint
\newif\ifPTMPOtags
\newif\ifPTMPOchi
\newif\ifPTMPOtimes
\newif\ifPTMPOarrow
\newif\ifPTMPOfinal
\PTMPOjointfalse
\PTMPOtagstrue
\PTMPOchitrue
\PTMPOtimestrue
\PTMPOarrowtrue
\PTMPOfinalfalse

\makeatletter
\define@key{ptmpo}{n}{\def\PTMPOn{#1}}
\define@key{ptmpo}{at}{\def\PTMPOat{#1}}
\define@key{ptmpo}{name}{\def\PTMPOname{#1}}
\define@key{ptmpo}{spacing}{\setlength{\PTMPOspacing}{#1}}
\define@key{ptmpo}{joint}{\csname PTMPOjoint#1\endcsname}
\define@key{ptmpo}{tags}{\csname PTMPOtags#1\endcsname}
\define@key{ptmpo}{chi}{\csname PTMPOchi#1\endcsname}
\define@key{ptmpo}{times}{\csname PTMPOtimes#1\endcsname}
\define@key{ptmpo}{time arrow}{\csname PTMPOarrow#1\endcsname}
\define@key{ptmpo}{final}{\csname PTMPOfinal#1\endcsname}
\define@key{ptmpo}{label}{\def\PTMPOlabel{#1}}
\makeatother

% Draw n time cores W^{[0]} (right, t_0) ... W^{[n-1]} (left).
% Node names:  <name>-W-<k>
% IO tips:     <name>-out-<k>  (process output, RIGHT of pair)
%              <name>-in-<k>   (process input,  LEFT  of pair)
% Time arrow:  <name>-t0  (right)  -->  <name>-tn (left)
\newcommand{\PTDrawProcessMPO}[1][]{%
  \PTMPOjointfalse
  \PTMPOtagstrue
  \PTMPOchitrue
  \PTMPOtimestrue
  \PTMPOarrowtrue
  \PTMPOfinalfalse
  \setlength{\PTMPOspacing}{\PTTimeSpacing}%
  \def\PTMPOn{3}%
  \def\PTMPOat{0,0}%
  \def\PTMPOname{pt}%
  \def\PTMPOlabel{}%
  \setkeys{ptmpo}{#1}%
  \coordinate (\PTMPOname-origin) at (\PTMPOat);
  \edef\PTMPOlast{\PTMPOname-W-\the\numexpr\PTMPOn-1\relax}%
  % --- cores, right (k=0) to left ------------------------------------------
  \foreach \k in {0,...,\the\numexpr\PTMPOn-1\relax} {%
    \path (\PTMPOname-origin)
      ++(-\k*\PTMPOspacing,0)
      coordinate (\PTMPOname-c-\k);
    \PTProcessCore{\PTMPOname-W-\k}{(\PTMPOname-c-\k)}{\(W^{[\k]}\)}%
    % Local ports: output on the right, input on the left (local time).
    \coordinate (\PTMPOname-outsrc-\k) at
      ([xshift=\PTIOSep]\PTMPOname-W-\k.south);
    \coordinate (\PTMPOname-insrc-\k) at
      ([xshift=-\PTIOSep]\PTMPOname-W-\k.south);
    \draw[PTLiouvilleLeg]
      (\PTMPOname-outsrc-\k) -- ++(0,-\PTLegLength)
      coordinate (\PTMPOname-out-\k);
    \draw[PTLiouvilleLeg]
      (\PTMPOname-insrc-\k) -- ++(0,-\PTLegLength)
      coordinate (\PTMPOname-in-\k);
  }%
  % --- joint body (no exposed link indices) or MPO links -------------------
  \ifPTMPOjoint
    \node[
      fit=(\PTMPOname-W-0) (\PTMPOlast),
      inner xsep=3.6pt,
      inner ysep=4.6pt,
      rounded corners=5.5pt,
      fill=PTProcess!55,
      draw=black,
      line width=0.7pt,
    ] (\PTMPOname-body) {};
    \draw[
      draw=PTProcess!70!black,
      line width=3.4pt,
      line cap=round,
    ] (\PTMPOlast.center) -- (\PTMPOname-W-0.center);
  \else
    \ifnum\PTMPOn>1
      \foreach \k [
        evaluate=\k as \knext using int(\k+1)
      ] in {0,...,\the\numexpr\PTMPOn-2\relax} {%
        \draw[PTLinkLeg]
          (\PTMPOname-W-\k.west) -- (\PTMPOname-W-\knext.east)
          coordinate[midway] (\PTMPOname-link-\k);
        \ifPTMPOchi
          \node[font=\small, text=PTLabel, inner sep=0.5pt, above=1.4pt]
            at (\PTMPOname-link-\k) {\(\chi_{\k}\)};
        \fi
        \ifPTMPOtags
          \PTIndexLabel[
            anchor=south,
            tag={PTLink,t-0\k},
            yshift=11.0pt,
          ]{\PTMPOname-link-\k}%
        \fi
      }%
    \fi
    % Dimension-1 boundary stubs.
    \draw[PTLinkLeg] (\PTMPOname-W-0.east) -- ++(4.2mm,0)
      coordinate (\PTMPOname-bnd-R);
    \begin{scope}[on PTforeground layer]
      \node[PTDummyBond] (\PTMPOname-dummy-R) at (\PTMPOname-bnd-R) {};
    \end{scope}
    \draw[PTLinkLeg]
      (\PTMPOlast.west) -- ++(-4.2mm,0)
      coordinate (\PTMPOname-bnd-L);
    \begin{scope}[on PTforeground layer]
      \node[PTDummyBond] (\PTMPOname-dummy-L) at (\PTMPOname-bnd-L) {};
    \end{scope}
  \fi
  % --- optional final (t_n) physical output on the far left ----------------
  \ifPTMPOfinal
    \coordinate (\PTMPOname-final-src) at ([yshift=1.6mm]\PTMPOlast.west);
    \draw[PTLiouvilleLeg] (\PTMPOname-final-src) -- ++(-7.5mm,0)
      coordinate (\PTMPOname-final);
  \fi
  % --- time labels and left-pointing arrow ---------------------------------
  \coordinate (\PTMPOname-t0) at ([yshift=-15.2mm, xshift=5.0mm]\PTMPOname-W-0.south);
  \coordinate (\PTMPOname-tn) at ([yshift=-15.2mm, xshift=-5.0mm]\PTMPOlast.south);
  \ifPTMPOtimes
    \node[font=\small, text=PTLabel, inner sep=0.4pt, anchor=west]
      at (\PTMPOname-t0) {\(t_{0}\)};
    \node[font=\small, text=PTLabel, inner sep=0.4pt, anchor=east]
      at (\PTMPOname-tn) {\(t_{n}\)};
  \fi
  \ifPTMPOarrow
    \PTTimeArrow[below=1.2pt]{\PTMPOname-t0}{\PTMPOname-tn}{\(t\)}%
  \fi
  \ifPTMPOtags
    \PTIndexLabel[
      anchor=west,
      math={\ell_{\mathrm{out}}},
      tag={Output,t-00},
      xshift=3.0pt,
    ]{\PTMPOname-out-0}%
    \PTIndexLabel[
      anchor=east,
      math={\ell_{\mathrm{in}}},
      tag={Input,t-00},
      xshift=-3.0pt,
      yshift=-4.2pt,
    ]{\PTMPOname-in-0}%
  \fi
  \def\PTMPOlabelEmpty{}%
  \ifx\PTMPOlabel\PTMPOlabelEmpty
  \else
    \node[font=\small, text=PTLabel, inner sep=1pt, above=3.6pt]
      at ($ (\PTMPOname-W-0.north)!0.5!(\PTMPOlast.north) $)
      {\PTMPOlabel};
  \fi
}

\newcommand{\PTDrawJointProcess}[1][]{%
  \PTDrawProcessMPO[joint=true,chi=false,tags=false,#1]%
}

% -----------------------------------------------------------------------------
% Example of a future instrument (commented; copy and adapt).
% -----------------------------------------------------------------------------
% \newcommand{\PTReset}[4][]{%
%   \PTInstrumentNode[fill=PTInstrument, classical=none, #1]{#2}{#3}{#4}%
% }
% % Tester / ancilla interaction (two-site box):
% \newcommand{\PTTester}[4][]{%
%   \PTInstrumentNode[
%     fill=PTInstrument,
%     minimum width=1.55\PTCoreWidth,
%     #1
%   ]{#2}{#3}{#4}%
% }
% % Feedback-conditioned operation: reuse classical=down and wire it into
% % a later \PTInstrumentNode.
%
% This file is a library.  Compile pt_tikz_demo.tex for a catalogue of
% every primitive.

% Figure fragments then:
%   \input{figures/tikz/<name>.tikz}
%
% =============================================================================

\ifdefined\PTTikZLoaded
  \expandafter\endinput
\fi
\def\PTTikZLoaded{1}

\RequirePackage{tikz}
\RequirePackage{xcolor}
\RequirePackage{keyval}
\usetikzlibrary{
  arrows.meta,
  backgrounds,
  calc,
  decorations.markings,
  fit,
  matrix,
  positioning,
  shapes.geometric
}

% -----------------------------------------------------------------------------
% Colours (Okabe--Ito / colour-blind friendly; intelligible in grayscale)
% -----------------------------------------------------------------------------
\definecolor{PTHilbertKet}{HTML}{0072B2}
\definecolor{PTHilbertBra}{HTML}{D55E00}
\definecolor{PTLiouville}{HTML}{7B3294}
\definecolor{PTLink}{HTML}{009E73}
\definecolor{PTTensor}{HTML}{E8E8E8}
\definecolor{PTProcess}{HTML}{DCEAF7}
\definecolor{PTInstrument}{HTML}{F6E8A6}
\definecolor{PTTrace}{HTML}{CC79A7}
\definecolor{PTMeasurement}{HTML}{E69F00}
\definecolor{PTLabel}{HTML}{333333}
\definecolor{PTBond}{HTML}{6A6A6A}

% -----------------------------------------------------------------------------
% Lengths (the single place to retune the visual language)
% -----------------------------------------------------------------------------
\newlength{\PTCoreWidth}      \setlength{\PTCoreWidth}{9.6mm}
\newlength{\PTCoreHeight}     \setlength{\PTCoreHeight}{9.6mm}
\newlength{\PTTimeSpacing}    \setlength{\PTTimeSpacing}{20.0mm}
\newlength{\PTLegLength}      \setlength{\PTLegLength}{7.0mm}
\newlength{\PTIOSep}          \setlength{\PTIOSep}{2.4mm}
% Equilateral triangle: vertical base = side, horizontal altitude = (sqrt(3)/2)*side.
\newlength{\PTKetHeight}      \setlength{\PTKetHeight}{7.8mm}
\newlength{\PTKetWidth}       \setlength{\PTKetWidth}{6.75mm}
\newlength{\PTBendDepth}      \setlength{\PTBendDepth}{5.5mm}
\newlength{\PTCombinerSize}   \setlength{\PTCombinerSize}{7.2mm}
\newlength{\PTHilbertLW}      \setlength{\PTHilbertLW}{1.35pt}
\newlength{\PTLiouvilleLW}    \setlength{\PTLiouvilleLW}{0.70pt}
\newlength{\PTLiouvilleGap}   \setlength{\PTLiouvilleGap}{1.55pt}
\newlength{\PTLinkLW}         \setlength{\PTLinkLW}{2.70pt}
\newlength{\PTBondLW}         \setlength{\PTBondLW}{1.35pt}
\newlength{\PTClassicalLW}    \setlength{\PTClassicalLW}{1.15pt}
\newlength{\PTCoreStroke}     \setlength{\PTCoreStroke}{1.45pt}
\newlength{\PTCoreRadius}     \setlength{\PTCoreRadius}{1.3pt}
\newlength{\PTBridgeHalo}     \setlength{\PTBridgeHalo}{4.2pt}

% Scratch length for the MPO drawer (not a public tuning knob).
\newlength{\PTMPOspacing}

% -----------------------------------------------------------------------------
% Layers: wires behind, labels in the middle, filled shapes in front.
% -----------------------------------------------------------------------------
\pgfdeclarelayer{PTforeground}
\pgfsetlayers{background,main,PTforeground}
\tikzset{
  on PTforeground layer/.style={
    execute at begin scope={\pgfonlayer{PTforeground}},
    execute at end scope={\endpgfonlayer},
  },
}

% -----------------------------------------------------------------------------
% Line and node styles
% -----------------------------------------------------------------------------
\tikzset{
  PTPicture/.style={
    baseline,
    line join=round,
    >={Stealth[length=2.0mm, width=1.45mm]},
  },
  % --- wires (black; structure, not colour, distinguishes spaces) ----------
  PTWire/.style={
    draw=black,
    line width=\PTHilbertLW,
    line cap=butt,
    line join=round,
    double=none,
  },
  PTHilbertKetLeg/.style={
    PTWire,
    on background layer,
  },
  PTHilbertBraLeg/.style={
    draw=black,
    line width=\PTHilbertLW,
    line cap=butt,
    line join=round,
    double=none,
    on background layer,
  },
  % Optional prime tick at the midpoint of a bra wire.
  PTHilbertBraLeg primed/.style={
    PTHilbertBraLeg,
    postaction={
      decorate,
      decoration={
        markings,
        mark=at position 0.55 with {
          \node[font=\scriptsize, text=PTLabel, inner sep=0.4pt,
                xshift=3.2pt, yshift=0.4pt] {\(\prime\)};
        }
      }
    },
  },
  PTLiouvilleLeg/.style={
    draw=black,
    line width=\PTLiouvilleLW,
    double=white,
    double distance=\PTLiouvilleGap,
    line cap=butt,
    line join=round,
    on background layer,
  },
  PTBondLeg/.style={
    draw=black,
    line width=\PTBondLW,
    line cap=butt,
    line join=round,
    double=none,
    on background layer,
  },
  PTLinkLeg/.style={
    draw=black,
    line width=\PTLinkLW,
    line cap=butt,
    line join=round,
    double=none,
    on background layer,
  },
  PTClassicalLeg/.style={
    draw=black,
    line width=\PTClassicalLW,
    dash pattern=on 2.4pt off 1.6pt,
    line cap=butt,
    line join=round,
    double=none,
    on background layer,
  },
  PTTraceCup/.style={
    draw=PTTrace,
    line width=\PTHilbertLW,
    line cap=round,
    line join=round,
    double=none,
    on background layer,
  },
  % White halo so an over-crossing does not look like a contraction.
  PTBridge/.style={
    preaction={
      draw=white,
      line width=\PTBridgeHalo,
      double=none,
      line cap=butt,
    },
  },
  % --- cores ---------------------------------------------------------------
  PTCore/.style={
    draw=black,
    fill=PTTensor,
    rounded corners=\PTCoreRadius,
    minimum width=\PTCoreWidth,
    minimum height=\PTCoreHeight,
    inner sep=2.2pt,
    outer sep=0pt,
    line width=\PTCoreStroke,
    font=\small,
    text=PTLabel,
    align=center,
  },
  PTTraceNode/.style={
    draw=black,
    fill=PTTrace!18,
    rounded corners=\PTCoreRadius,
    minimum width=7.2mm,
    minimum height=7.2mm,
    inner sep=1.4pt,
    outer sep=0pt,
    line width=\PTCoreStroke,
    font=\small,
    text=PTLabel,
    align=center,
  },
  PTDummyBond/.style={
    circle,
    fill=black,
    draw=none,
    inner sep=1.15pt,
    outer sep=0pt,
  },
}

% -----------------------------------------------------------------------------
% Index labels: math + ITensor tag + optional prime level
% -----------------------------------------------------------------------------
\makeatletter
\define@key{ptindex}{anchor}{\def\PTIndexAnchor{#1}}
\define@key{ptindex}{math}{\def\PTIndexMath{#1}}
\define@key{ptindex}{tag}{\def\PTIndexTag{#1}}
\define@key{ptindex}{prime}{\def\PTIndexPrime{#1}}
\define@key{ptindex}{xshift}{\def\PTIndexX{#1}}
\define@key{ptindex}{yshift}{\def\PTIndexY{#1}}
\makeatother

% Append a visual prime level to math (') and to a tag (ASCII apostrophe).
\newcommand{\PTIndexPrimeTokens}{%
  \ifx\PTIndexPrime\empty
  \else
    \ifnum\PTIndexPrime>0 '\fi
    \ifnum\PTIndexPrime>1 '\fi
    \ifnum\PTIndexPrime>2 '\fi
  \fi
}

\newcommand{\PTIndexLabel}[2][]{%
  \def\PTIndexAnchor{south}%
  \def\PTIndexMath{}%
  \def\PTIndexTag{}%
  \def\PTIndexPrime{}%
  \def\PTIndexX{0pt}%
  \def\PTIndexY{0pt}%
  \setkeys{ptindex}{#1}%
  \node[
    inner sep=0.9pt,
    outer sep=1.2pt,
    anchor=\PTIndexAnchor,
    align=center,
    xshift=\PTIndexX,
    yshift=\PTIndexY,
    text=PTLabel,
  ] at (#2) {%
    \def\PTIndexMathEmpty{}%
    \ifx\PTIndexMath\PTIndexMathEmpty
    \else
      {\small\(\PTIndexMath\PTIndexPrimeTokens\)}%
      \ifx\PTIndexTag\PTIndexMathEmpty
      \else \\[-0.15ex]%
      \fi
    \fi
    \ifx\PTIndexTag\PTIndexMathEmpty
    \else
      {\scriptsize\ttfamily \PTIndexTag\PTIndexPrimeTokens}%
    \fi
  };%
}

% -----------------------------------------------------------------------------
% Ket and bra boundary tensors
% -----------------------------------------------------------------------------
% Invisible rectangular node provides named anchors.  The triangle is drawn
% on top so that .west (ket) / .east (bra) is the wire attachment (base).

\newcommand{\PTKet}[3]{%
  \node[
    name=#1,
    inner sep=0pt,
    outer sep=0pt,
    minimum width=\PTKetWidth,
    minimum height=\PTKetHeight,
    font=\small,
  ] at #2 {};
  \begin{scope}[on PTforeground layer]
    \filldraw[
      draw=black,
      fill=PTHilbertKet!18,
      line width=\PTCoreStroke,
      line join=round,
    ] (#1.west |- #1.south) -- (#1.east) -- (#1.west |- #1.north) -- cycle;
    \node[font=\small, text=PTLabel, inner sep=0.4pt]
      at ($ (#1.west)!0.36!(#1.east) $) {#3};
  \end{scope}
}

\newcommand{\PTBra}[3]{%
  \node[
    name=#1,
    inner sep=0pt,
    outer sep=0pt,
    minimum width=\PTKetWidth,
    minimum height=\PTKetHeight,
    font=\small,
  ] at #2 {};
  \begin{scope}[on PTforeground layer]
    \filldraw[
      draw=black,
      fill=PTHilbertBra!18,
      line width=\PTCoreStroke,
      line join=round,
    ] (#1.east |- #1.south) -- (#1.west) -- (#1.east |- #1.north) -- cycle;
    \node[font=\small, text=PTLabel, inner sep=0.4pt]
      at ($ (#1.east)!0.36!(#1.west) $) {#3};
  \end{scope}
}

\newcommand{\PTLiouvilleKet}[3]{%
  \node[
    name=#1,
    inner sep=0pt,
    outer sep=0pt,
    minimum width=\PTKetWidth,
    minimum height=\PTKetHeight,
    font=\small,
  ] at #2 {};
  \begin{scope}[on PTforeground layer]
    \filldraw[
      draw=black,
      fill=PTLiouville!16,
      line width=\PTCoreStroke,
      line join=round,
    ] (#1.west |- #1.south) -- (#1.east) -- (#1.west |- #1.north) -- cycle;
    \node[font=\small, text=PTLabel, inner sep=0.4pt]
      at ($ (#1.west)!0.36!(#1.east) $) {#3};
  \end{scope}
}

\newcommand{\PTLiouvilleBra}[3]{%
  \node[
    name=#1,
    inner sep=0pt,
    outer sep=0pt,
    minimum width=\PTKetWidth,
    minimum height=\PTKetHeight,
    font=\small,
  ] at #2 {};
  \begin{scope}[on PTforeground layer]
    \filldraw[
      draw=black,
      fill=PTLiouville!16,
      line width=\PTCoreStroke,
      line join=round,
    ] (#1.east |- #1.south) -- (#1.west) -- (#1.east |- #1.north) -- cycle;
    \node[font=\small, text=PTLabel, inner sep=0.4pt]
      at ($ (#1.east)!0.36!(#1.west) $) {#3};
  \end{scope}
}

% -----------------------------------------------------------------------------
% Square tensor cores (thick stroke, slight corner rounding)
% -----------------------------------------------------------------------------
\makeatletter
\define@key{pttensor}{fill}{\def\PTTensorFill{#1}}
\define@key{pttensor}{draw}{\def\PTTensorDrawCol{#1}}
\define@key{pttensor}{minimum width}{\def\PTTensorMinW{#1}}
\define@key{pttensor}{minimum height}{\def\PTTensorMinH{#1}}
\define@key{pttensor}{rounded}{\def\PTTensorRnd{#1}}
\define@key{pttensor}{font}{\def\PTTensorFont{#1}}
\makeatother

\newcommand{\PTTensor}[4][]{%
  \def\PTTensorFill{PTTensor}%
  \def\PTTensorDrawCol{black}%
  \def\PTTensorMinW{\PTCoreWidth}%
  \def\PTTensorMinH{\PTCoreHeight}%
  \def\PTTensorRnd{\PTCoreRadius}%
  \def\PTTensorFont{\small}%
  \setkeys{pttensor}{#1}%
  \begin{scope}[on PTforeground layer]
    \node[
      name=#2,
      PTCore,
      fill=\PTTensorFill,
      draw=\PTTensorDrawCol,
      minimum width=\PTTensorMinW,
      minimum height=\PTTensorMinH,
      rounded corners=\PTTensorRnd,
      font=\PTTensorFont,
    ] at #3 {#4};
  \end{scope}
}

\newcommand{\PTOperator}[4][]{%
  \PTTensor[fill=PTTensor, minimum width=\PTCoreWidth, minimum height=\PTCoreWidth, #1]{#2}{#3}{#4}%
}

\newcommand{\PTSuperoperator}[4][]{%
  \PTTensor[
    fill=PTTensor,
    minimum width=\PTCoreWidth,
    minimum height=\PTCoreWidth,
    #1
  ]{#2}{#3}{#4}%
}

\newcommand{\PTProcessCore}[4][]{%
  \PTTensor[
    fill=PTProcess,
    minimum width=\PTCoreWidth,
    minimum height=\PTCoreWidth,
    #1
  ]{#2}{#3}{#4}%
}

\newcommand{\PTInstrument}[4][]{%
  \PTTensor[
    fill=PTInstrument,
    minimum width=\PTCoreWidth,
    minimum height=\PTCoreWidth,
    #1
  ]{#2}{#3}{#4}%
}

% Vertical Hilbert ket (north) and bra (south) stubs on an operator.
\newcommand{\PTOperatorLegs}[1]{%
  \draw[PTHilbertKetLeg] (#1.north) -- ++(0,\PTLegLength) coordinate (#1-ket);
  \draw[PTHilbertBraLeg] (#1.south) -- ++(0,-\PTLegLength) coordinate (#1-bra);
}

% -----------------------------------------------------------------------------
% Combiners: s \otimes s'  -->  \ell    and the inverse
% -----------------------------------------------------------------------------
\makeatletter
\define@key{ptcomb}{direction}{\def\PTCombDir{#1}}
\makeatother
\def\PTCombDirUp{up}
\def\PTCombDirDown{down}
\def\PTCombDirLeft{left}
\def\PTCombDirRight{right}

\newcommand{\PTCombiner}[4][]{%
  \def\PTCombDir{up}%
  \setkeys{ptcomb}{#1}%
  \ifx\PTCombDir\PTCombDirDown
    \def\PTCombRot{180}%
  \else\ifx\PTCombDir\PTCombDirLeft
    \def\PTCombRot{90}%
  \else\ifx\PTCombDir\PTCombDirRight
    \def\PTCombRot{270}%
  \else
    \def\PTCombRot{0}%
  \fi\fi\fi
  \begin{scope}[on PTforeground layer]
    \node[
      name=#2,
      isosceles triangle,
      isosceles triangle apex angle=60,
      shape border rotate=\PTCombRot,
      inner sep=0pt,
      outer sep=0pt,
      minimum width=\PTCombinerSize,
      draw=black,
      fill=white,
      line width=\PTCoreStroke,
      font=\scriptsize,
      text=PTLabel,
    ] at #3 {};
    \node[font=\scriptsize, text=PTLabel, inner sep=0.8pt, anchor=east]
      at ([xshift=-1.6pt]#2.west) {#4};
  \end{scope}
  \coordinate (#2-fused) at (#2.apex);
  \coordinate (#2-ket)   at (#2.left corner);
  \coordinate (#2-bra)   at (#2.right corner);
}

\newcommand{\PTCombinerInv}[4][]{%
  \PTCombiner[direction=down, #1]{#2}{#3}{#4}%
}

% -----------------------------------------------------------------------------
% Cups, caps, time arrow, callouts
% -----------------------------------------------------------------------------
% Cup (union, opens upward): used for traces and vectorisation bends.
% The two endpoints are the open wire tips; control points sit below them.
\newcommand{\PTCup}[3][]{%
  \draw[#1] (#2)
    .. controls ($ (#2) + (0,-\PTBendDepth) $)
            and ($ (#3) + (0,-\PTBendDepth) $) ..
    (#3);
}

% Trace: left and right semicircles joined by a straight bar BELOW the core
% (never behind it).  #2-loopmid is the bar midpoint for a single index tag.
% Drawn on the main layer so the bar is not covered by the core; the core
% still sits in front at the ports.
\newcommand{\PTTraceLoop}[2][PTWire]{%
  \draw[#1]
    (#2.west)
    arc[start angle=90, delta angle=180, radius=4.0mm]
    -- ([yshift=-8.0mm]#2.east)
    arc[start angle=270, delta angle=180, radius=4.0mm];
  \coordinate (#2-loopmid) at ([yshift=-8.0mm]#2.center);
}

% Transpose bends: left (ket) port over the top towards the right, now a
% primed bra (coordinate #1-i); right (bra) port under the bottom towards
% the left, now an unprimed ket (coordinate #1-j).
\newcommand{\PTTransposeBends}[1]{%
  \draw[PTHilbertBraLeg]
    (#1.west)
    arc[start angle=270, delta angle=-90, radius=3.3mm]
    arc[start angle=180, delta angle=-90, radius=3.3mm]
    coordinate (#1-itop);
  \draw[PTHilbertBraLeg primed]
    (#1-itop) -- ++(11.5mm,0) coordinate (#1-i);
  \draw[PTHilbertKetLeg]
    (#1.east)
    arc[start angle=90, delta angle=-90, radius=3.3mm]
    arc[start angle=0, delta angle=-90, radius=3.3mm]
    -- ++(-11.5mm,0) coordinate (#1-j);
}
\newcommand{\PTCap}[3][]{%
  \draw[#1] (#2)
    .. controls ($ (#2) + (0,\PTBendDepth) $)
            and ($ (#3) + (0,\PTBendDepth) $) ..
    (#3);
}

% Sideways cup used when bending a vertical bra wire into a second output.
% #2 = start (usually the operator's south / bra port).
% #3 = end   (the new upward ket-like port for the bra index).
% The bend goes to the LEFT of both points by default (column-major: |k>
% becomes the left factor).
\newcommand{\PTVecBend}[3][]{%
  \draw[#1] (#2)
    .. controls ($ (#2) + (0,-\PTBendDepth) $)
            and ($ (#3) + (0,-\PTBendDepth) $) ..
    (#3);
}

% Time arrow: ALWAYS left-pointing.  #1 = right coord, #2 = left coord.
% Optional first argument is a TikZ node-anchor for the label (default below).
\newcommand{\PTTimeArrow}[4][below=2.4pt]{%
  \draw[
    -{Stealth[length=2.2mm, width=1.55mm]},
    line width=1.05pt,
    draw=PTLabel,
    line cap=round,
  ] (#2) -- (#3)
    node[midway, #1, font=\small, text=PTLabel, inner sep=0.6pt] {#4};
}

\newcommand{\PTPanelLabel}[2]{%
  \node[
    font=\small\bfseries,
    text=PTLabel,
    inner sep=0.4pt,
    anchor=south west,
  ] at #1 {#2};
}

\newcommand{\PTCallout}[3][]{%
  \node[
    font=\scriptsize,
    text=PTLabel,
    inner sep=0.8pt,
    align=left,
    #1,
  ] at (#2) {#3};
}

\newcommand{\PTMapsTo}[2][]{%
  \node[font=\small, text=PTLabel, inner sep=1pt, #1] at #2 {\(\longmapsto\)};
}

% -----------------------------------------------------------------------------
% Instruments
% -----------------------------------------------------------------------------
\newif\ifPTInstClassical
\PTInstClassicalfalse
\def\PTInstKindBox{box}
\def\PTInstKindKet{ket}
\def\PTInstKindBra{bra}
\def\PTInstKindTrace{trace}

\makeatletter
\define@key{ptinst}{fill}{\def\PTInstFill{#1}}
\define@key{ptinst}{draw}{\def\PTInstDraw{#1}}
\define@key{ptinst}{minimum width}{\def\PTInstMinW{#1}}
\define@key{ptinst}{minimum height}{\def\PTInstMinH{#1}}
\define@key{ptinst}{kind}{\def\PTInstKind{#1}}
\define@key{ptinst}{classical}{\def\PTInstClassDir{#1}}
\define@key{ptinst}{classical label}{\def\PTInstClassLab{#1}}
\define@key{ptinst}{classical length}{\def\PTInstClassLen{#1}}
\makeatother

% Generic constructor.  Quantum legs are NOT auto-drawn: connect to the
% node anchors (north, south, east, west, or the triangle bases).
% Classical legs ARE auto-drawn when classical != none.
\newcommand{\PTInstrumentNode}[4][]{%
  \def\PTInstFill{PTInstrument}%
  \def\PTInstDraw{black}%
  \def\PTInstMinW{\PTCoreWidth}%
  \def\PTInstMinH{\PTCoreWidth}%
  \def\PTInstKind{box}%
  \def\PTInstClassDir{none}%
  \def\PTInstClassLab{}%
  \def\PTInstClassLen{6.2mm}%
  \setkeys{ptinst}{#1}%
  \def\PTInstNone{none}%
  \ifx\PTInstKind\PTInstKindKet
    \node[
      name=#2,
      inner sep=0pt,
      outer sep=0pt,
      minimum width=\PTKetWidth,
      minimum height=\PTKetHeight,
    ] at #3 {};
    \begin{scope}[on PTforeground layer]
      \filldraw[
        draw=black,
        fill=\PTInstFill,
        line width=\PTCoreStroke,
        line join=round,
      ] (#2.west |- #2.south) -- (#2.east) -- (#2.west |- #2.north) -- cycle;
      \node[font=\small, text=PTLabel, inner sep=0.4pt]
        at ($ (#2.west)!0.36!(#2.east) $) {#4};
    \end{scope}
  \else
    \ifx\PTInstKind\PTInstKindBra
      \node[
        name=#2,
        inner sep=0pt,
        outer sep=0pt,
        minimum width=\PTKetWidth,
        minimum height=\PTKetHeight,
      ] at #3 {};
      \begin{scope}[on PTforeground layer]
        \filldraw[
          draw=black,
          fill=\PTInstFill,
          line width=\PTCoreStroke,
          line join=round,
        ] (#2.east |- #2.south) -- (#2.west) -- (#2.east |- #2.north) -- cycle;
        \node[font=\small, text=PTLabel, inner sep=0.4pt]
          at ($ (#2.east)!0.36!(#2.west) $) {#4};
      \end{scope}
    \else
      \ifx\PTInstKind\PTInstKindTrace
        \begin{scope}[on PTforeground layer]
          \node[name=#2, PTTraceNode] at #3 {#4};
        \end{scope}
      \else
        \begin{scope}[on PTforeground layer]
          \node[
            name=#2,
            PTCore,
            fill=\PTInstFill,
            draw=\PTInstDraw,
            minimum width=\PTInstMinW,
            minimum height=\PTInstMinH,
          ] at #3 {#4};
        \end{scope}
      \fi
    \fi
  \fi
  \ifx\PTInstClassDir\PTInstNone
  \else
    \def\PTInstDirDown{down}%
    \def\PTInstDirUp{up}%
    \def\PTInstDirLeft{left}%
    \def\PTInstDirRight{right}%
    \ifx\PTInstClassDir\PTInstDirDown
      \draw[PTClassicalLeg] (#2.south) -- ++(0,-\PTInstClassLen)
        coordinate (#2-classical);
    \else\ifx\PTInstClassDir\PTInstDirUp
      \draw[PTClassicalLeg] (#2.north) -- ++(0,\PTInstClassLen)
        coordinate (#2-classical);
    \else\ifx\PTInstClassDir\PTInstDirLeft
      \draw[PTClassicalLeg] (#2.west) -- ++(-\PTInstClassLen,0)
        coordinate (#2-classical);
    \else
      \draw[PTClassicalLeg] (#2.east) -- ++(\PTInstClassLen,0)
        coordinate (#2-classical);
    \fi\fi\fi
    \def\PTInstClassEmpty{}%
    \ifx\PTInstClassLab\PTInstClassEmpty
    \else
      \node[font=\small, text=PTLabel, inner sep=0.8pt, anchor=north]
        at ([yshift=-1.0pt]#2-classical) {\(\PTInstClassLab\)};
    \fi
  \fi
}

\newcommand{\PTIdentity}[4][]{%
  \PTInstrumentNode[fill=PTTensor, #1]{#2}{#3}{#4}%
}

\newcommand{\PTControl}[4][]{%
  \PTInstrumentNode[#1]{#2}{#3}{#4}%
}

\newcommand{\PTUnitary}[4][]{%
  \PTInstrumentNode[#1]{#2}{#3}{#4}%
}

\newcommand{\PTPrepare}[4][]{%
  \PTInstrumentNode[kind=ket, fill=PTInstrument, #1]{#2}{#3}{#4}%
}

\newcommand{\PTMeasure}[4][]{%
  \PTInstrumentNode[
    classical=down,
    classical label={x},
    #1
  ]{#2}{#3}{#4}%
}

\newcommand{\PTProjective}[4][]{%
  \PTInstrumentNode[
    classical=down,
    classical label={x},
    #1
  ]{#2}{#3}{#4}%
}

% Measurement (earlier, right) then preparation (later, left).
% Quantum wire is broken; the classical outcome is drawn separately.
% Nodes: <name>-E, <name>-P, <name>-x.
\newcommand{\PTCausalBreak}[3][]{%
  \begin{scope}[shift={#3}]
    \PTInstrumentNode[
      classical=none,
      minimum width=8.4mm,
      minimum height=8.4mm,
      #1
    ]{#2-E}{(0,0)}{\(E^{(x)}\)}
    \PTPrepare{#2-P}{(-1.65,0)}{\(\mathcal{P}\)}
    \draw[PTClassicalLeg] (#2-E.south) -- ++(0,-5.8mm) coordinate (#2-x);
    \node[font=\small, text=PTLabel, inner sep=0.6pt, anchor=north]
      at ([yshift=-0.8pt]#2-x) {\(x\)};
  \end{scope}
}

% Hilbert-space trace: magenta cup closing ket and bra tips.
\newcommand{\PTTraceHilbert}[2]{%
  \PTCup[PTTraceCup]{#1}{#2}%
  \node[font=\small, text=PTTrace, inner sep=0.5pt, anchor=north]
    at ($ (#1)!0.5!(#2) + (0,-\PTBendDepth-0.5pt) $) {\(\mathrm{Tr}\)};
}

% Liouville-space trace: compact Tr node sitting on a fused leg.
\newcommand{\PTTraceLiouville}[3]{%
  \begin{scope}[on PTforeground layer]
    \node[name=#1, PTTraceNode] at #2 {#3};
  \end{scope}
}

% -----------------------------------------------------------------------------
% PT-MPO / joint process tensor
% -----------------------------------------------------------------------------
\newif\ifPTMPOjoint
\newif\ifPTMPOtags
\newif\ifPTMPOchi
\newif\ifPTMPOtimes
\newif\ifPTMPOarrow
\newif\ifPTMPOfinal
\PTMPOjointfalse
\PTMPOtagstrue
\PTMPOchitrue
\PTMPOtimestrue
\PTMPOarrowtrue
\PTMPOfinalfalse

\makeatletter
\define@key{ptmpo}{n}{\def\PTMPOn{#1}}
\define@key{ptmpo}{at}{\def\PTMPOat{#1}}
\define@key{ptmpo}{name}{\def\PTMPOname{#1}}
\define@key{ptmpo}{spacing}{\setlength{\PTMPOspacing}{#1}}
\define@key{ptmpo}{joint}{\csname PTMPOjoint#1\endcsname}
\define@key{ptmpo}{tags}{\csname PTMPOtags#1\endcsname}
\define@key{ptmpo}{chi}{\csname PTMPOchi#1\endcsname}
\define@key{ptmpo}{times}{\csname PTMPOtimes#1\endcsname}
\define@key{ptmpo}{time arrow}{\csname PTMPOarrow#1\endcsname}
\define@key{ptmpo}{final}{\csname PTMPOfinal#1\endcsname}
\define@key{ptmpo}{label}{\def\PTMPOlabel{#1}}
\makeatother

% Draw n time cores W^{[0]} (right, t_0) ... W^{[n-1]} (left).
% Node names:  <name>-W-<k>
% IO tips:     <name>-out-<k>  (process output, RIGHT of pair)
%              <name>-in-<k>   (process input,  LEFT  of pair)
% Time arrow:  <name>-t0  (right)  -->  <name>-tn (left)
\newcommand{\PTDrawProcessMPO}[1][]{%
  \PTMPOjointfalse
  \PTMPOtagstrue
  \PTMPOchitrue
  \PTMPOtimestrue
  \PTMPOarrowtrue
  \PTMPOfinalfalse
  \setlength{\PTMPOspacing}{\PTTimeSpacing}%
  \def\PTMPOn{3}%
  \def\PTMPOat{0,0}%
  \def\PTMPOname{pt}%
  \def\PTMPOlabel{}%
  \setkeys{ptmpo}{#1}%
  \coordinate (\PTMPOname-origin) at (\PTMPOat);
  \edef\PTMPOlast{\PTMPOname-W-\the\numexpr\PTMPOn-1\relax}%
  % --- cores, right (k=0) to left ------------------------------------------
  \foreach \k in {0,...,\the\numexpr\PTMPOn-1\relax} {%
    \path (\PTMPOname-origin)
      ++(-\k*\PTMPOspacing,0)
      coordinate (\PTMPOname-c-\k);
    \PTProcessCore{\PTMPOname-W-\k}{(\PTMPOname-c-\k)}{\(W^{[\k]}\)}%
    % Local ports: output on the right, input on the left (local time).
    \coordinate (\PTMPOname-outsrc-\k) at
      ([xshift=\PTIOSep]\PTMPOname-W-\k.south);
    \coordinate (\PTMPOname-insrc-\k) at
      ([xshift=-\PTIOSep]\PTMPOname-W-\k.south);
    \draw[PTLiouvilleLeg]
      (\PTMPOname-outsrc-\k) -- ++(0,-\PTLegLength)
      coordinate (\PTMPOname-out-\k);
    \draw[PTLiouvilleLeg]
      (\PTMPOname-insrc-\k) -- ++(0,-\PTLegLength)
      coordinate (\PTMPOname-in-\k);
  }%
  % --- joint body (no exposed link indices) or MPO links -------------------
  \ifPTMPOjoint
    \node[
      fit=(\PTMPOname-W-0) (\PTMPOlast),
      inner xsep=3.6pt,
      inner ysep=4.6pt,
      rounded corners=5.5pt,
      fill=PTProcess!55,
      draw=black,
      line width=0.7pt,
    ] (\PTMPOname-body) {};
    \draw[
      draw=PTProcess!70!black,
      line width=3.4pt,
      line cap=round,
    ] (\PTMPOlast.center) -- (\PTMPOname-W-0.center);
  \else
    \ifnum\PTMPOn>1
      \foreach \k [
        evaluate=\k as \knext using int(\k+1)
      ] in {0,...,\the\numexpr\PTMPOn-2\relax} {%
        \draw[PTLinkLeg]
          (\PTMPOname-W-\k.west) -- (\PTMPOname-W-\knext.east)
          coordinate[midway] (\PTMPOname-link-\k);
        \ifPTMPOchi
          \node[font=\small, text=PTLabel, inner sep=0.5pt, above=1.4pt]
            at (\PTMPOname-link-\k) {\(\chi_{\k}\)};
        \fi
        \ifPTMPOtags
          \PTIndexLabel[
            anchor=south,
            tag={PTLink,t-0\k},
            yshift=11.0pt,
          ]{\PTMPOname-link-\k}%
        \fi
      }%
    \fi
    % Dimension-1 boundary stubs.
    \draw[PTLinkLeg] (\PTMPOname-W-0.east) -- ++(4.2mm,0)
      coordinate (\PTMPOname-bnd-R);
    \begin{scope}[on PTforeground layer]
      \node[PTDummyBond] (\PTMPOname-dummy-R) at (\PTMPOname-bnd-R) {};
    \end{scope}
    \draw[PTLinkLeg]
      (\PTMPOlast.west) -- ++(-4.2mm,0)
      coordinate (\PTMPOname-bnd-L);
    \begin{scope}[on PTforeground layer]
      \node[PTDummyBond] (\PTMPOname-dummy-L) at (\PTMPOname-bnd-L) {};
    \end{scope}
  \fi
  % --- optional final (t_n) physical output on the far left ----------------
  \ifPTMPOfinal
    \coordinate (\PTMPOname-final-src) at ([yshift=1.6mm]\PTMPOlast.west);
    \draw[PTLiouvilleLeg] (\PTMPOname-final-src) -- ++(-7.5mm,0)
      coordinate (\PTMPOname-final);
  \fi
  % --- time labels and left-pointing arrow ---------------------------------
  \coordinate (\PTMPOname-t0) at ([yshift=-15.2mm, xshift=5.0mm]\PTMPOname-W-0.south);
  \coordinate (\PTMPOname-tn) at ([yshift=-15.2mm, xshift=-5.0mm]\PTMPOlast.south);
  \ifPTMPOtimes
    \node[font=\small, text=PTLabel, inner sep=0.4pt, anchor=west]
      at (\PTMPOname-t0) {\(t_{0}\)};
    \node[font=\small, text=PTLabel, inner sep=0.4pt, anchor=east]
      at (\PTMPOname-tn) {\(t_{n}\)};
  \fi
  \ifPTMPOarrow
    \PTTimeArrow[below=1.2pt]{\PTMPOname-t0}{\PTMPOname-tn}{\(t\)}%
  \fi
  \ifPTMPOtags
    \PTIndexLabel[
      anchor=west,
      math={\ell_{\mathrm{out}}},
      tag={Output,t-00},
      xshift=3.0pt,
    ]{\PTMPOname-out-0}%
    \PTIndexLabel[
      anchor=east,
      math={\ell_{\mathrm{in}}},
      tag={Input,t-00},
      xshift=-3.0pt,
      yshift=-4.2pt,
    ]{\PTMPOname-in-0}%
  \fi
  \def\PTMPOlabelEmpty{}%
  \ifx\PTMPOlabel\PTMPOlabelEmpty
  \else
    \node[font=\small, text=PTLabel, inner sep=1pt, above=3.6pt]
      at ($ (\PTMPOname-W-0.north)!0.5!(\PTMPOlast.north) $)
      {\PTMPOlabel};
  \fi
}

\newcommand{\PTDrawJointProcess}[1][]{%
  \PTDrawProcessMPO[joint=true,chi=false,tags=false,#1]%
}

% -----------------------------------------------------------------------------
% Example of a future instrument (commented; copy and adapt).
% -----------------------------------------------------------------------------
% \newcommand{\PTReset}[4][]{%
%   \PTInstrumentNode[fill=PTInstrument, classical=none, #1]{#2}{#3}{#4}%
% }
% % Tester / ancilla interaction (two-site box):
% \newcommand{\PTTester}[4][]{%
%   \PTInstrumentNode[
%     fill=PTInstrument,
%     minimum width=1.55\PTCoreWidth,
%     #1
%   ]{#2}{#3}{#4}%
% }
% % Feedback-conditioned operation: reuse classical=down and wire it into
% % a later \PTInstrumentNode.
%
% This file is a library.  Compile pt_tikz_demo.tex for a catalogue of
% every primitive.

% Figure fragments then:
%   \input{figures/tikz/<name>.tikz}
%
% =============================================================================

\ifdefined\PTTikZLoaded
  \expandafter\endinput
\fi
\def\PTTikZLoaded{1}

\RequirePackage{tikz}
\RequirePackage{xcolor}
\RequirePackage{keyval}
\usetikzlibrary{
  arrows.meta,
  backgrounds,
  calc,
  decorations.markings,
  fit,
  matrix,
  positioning,
  shapes.geometric
}

% -----------------------------------------------------------------------------
% Colours (Okabe--Ito / colour-blind friendly; intelligible in grayscale)
% -----------------------------------------------------------------------------
\definecolor{PTHilbertKet}{HTML}{0072B2}
\definecolor{PTHilbertBra}{HTML}{D55E00}
\definecolor{PTLiouville}{HTML}{7B3294}
\definecolor{PTLink}{HTML}{009E73}
\definecolor{PTTensor}{HTML}{E8E8E8}
\definecolor{PTProcess}{HTML}{DCEAF7}
\definecolor{PTInstrument}{HTML}{F6E8A6}
\definecolor{PTTrace}{HTML}{CC79A7}
\definecolor{PTMeasurement}{HTML}{E69F00}
\definecolor{PTLabel}{HTML}{333333}
\definecolor{PTBond}{HTML}{6A6A6A}

% -----------------------------------------------------------------------------
% Lengths (the single place to retune the visual language)
% -----------------------------------------------------------------------------
\newlength{\PTCoreWidth}      \setlength{\PTCoreWidth}{9.6mm}
\newlength{\PTCoreHeight}     \setlength{\PTCoreHeight}{9.6mm}
\newlength{\PTTimeSpacing}    \setlength{\PTTimeSpacing}{20.0mm}
\newlength{\PTLegLength}      \setlength{\PTLegLength}{7.0mm}
\newlength{\PTIOSep}          \setlength{\PTIOSep}{2.4mm}
% Equilateral triangle: vertical base = side, horizontal altitude = (sqrt(3)/2)*side.
\newlength{\PTKetHeight}      \setlength{\PTKetHeight}{7.8mm}
\newlength{\PTKetWidth}       \setlength{\PTKetWidth}{6.75mm}
\newlength{\PTBendDepth}      \setlength{\PTBendDepth}{5.5mm}
\newlength{\PTCombinerSize}   \setlength{\PTCombinerSize}{7.2mm}
\newlength{\PTHilbertLW}      \setlength{\PTHilbertLW}{1.35pt}
\newlength{\PTLiouvilleLW}    \setlength{\PTLiouvilleLW}{0.70pt}
\newlength{\PTLiouvilleGap}   \setlength{\PTLiouvilleGap}{1.55pt}
\newlength{\PTLinkLW}         \setlength{\PTLinkLW}{2.70pt}
\newlength{\PTBondLW}         \setlength{\PTBondLW}{1.35pt}
\newlength{\PTClassicalLW}    \setlength{\PTClassicalLW}{1.15pt}
\newlength{\PTCoreStroke}     \setlength{\PTCoreStroke}{1.45pt}
\newlength{\PTCoreRadius}     \setlength{\PTCoreRadius}{1.3pt}
\newlength{\PTBridgeHalo}     \setlength{\PTBridgeHalo}{4.2pt}

% Scratch length for the MPO drawer (not a public tuning knob).
\newlength{\PTMPOspacing}

% -----------------------------------------------------------------------------
% Layers: wires behind, labels in the middle, filled shapes in front.
% -----------------------------------------------------------------------------
\pgfdeclarelayer{PTforeground}
\pgfsetlayers{background,main,PTforeground}
\tikzset{
  on PTforeground layer/.style={
    execute at begin scope={\pgfonlayer{PTforeground}},
    execute at end scope={\endpgfonlayer},
  },
}

% -----------------------------------------------------------------------------
% Line and node styles
% -----------------------------------------------------------------------------
\tikzset{
  PTPicture/.style={
    baseline,
    line join=round,
    >={Stealth[length=2.0mm, width=1.45mm]},
  },
  % --- wires (black; structure, not colour, distinguishes spaces) ----------
  PTWire/.style={
    draw=black,
    line width=\PTHilbertLW,
    line cap=butt,
    line join=round,
    double=none,
  },
  PTHilbertKetLeg/.style={
    PTWire,
    on background layer,
  },
  PTHilbertBraLeg/.style={
    draw=black,
    line width=\PTHilbertLW,
    line cap=butt,
    line join=round,
    double=none,
    on background layer,
  },
  % Optional prime tick at the midpoint of a bra wire.
  PTHilbertBraLeg primed/.style={
    PTHilbertBraLeg,
    postaction={
      decorate,
      decoration={
        markings,
        mark=at position 0.55 with {
          \node[font=\scriptsize, text=PTLabel, inner sep=0.4pt,
                xshift=3.2pt, yshift=0.4pt] {\(\prime\)};
        }
      }
    },
  },
  PTLiouvilleLeg/.style={
    draw=black,
    line width=\PTLiouvilleLW,
    double=white,
    double distance=\PTLiouvilleGap,
    line cap=butt,
    line join=round,
    on background layer,
  },
  PTBondLeg/.style={
    draw=black,
    line width=\PTBondLW,
    line cap=butt,
    line join=round,
    double=none,
    on background layer,
  },
  PTLinkLeg/.style={
    draw=black,
    line width=\PTLinkLW,
    line cap=butt,
    line join=round,
    double=none,
    on background layer,
  },
  PTClassicalLeg/.style={
    draw=black,
    line width=\PTClassicalLW,
    dash pattern=on 2.4pt off 1.6pt,
    line cap=butt,
    line join=round,
    double=none,
    on background layer,
  },
  PTTraceCup/.style={
    draw=PTTrace,
    line width=\PTHilbertLW,
    line cap=round,
    line join=round,
    double=none,
    on background layer,
  },
  % White halo so an over-crossing does not look like a contraction.
  PTBridge/.style={
    preaction={
      draw=white,
      line width=\PTBridgeHalo,
      double=none,
      line cap=butt,
    },
  },
  % --- cores ---------------------------------------------------------------
  PTCore/.style={
    draw=black,
    fill=PTTensor,
    rounded corners=\PTCoreRadius,
    minimum width=\PTCoreWidth,
    minimum height=\PTCoreHeight,
    inner sep=2.2pt,
    outer sep=0pt,
    line width=\PTCoreStroke,
    font=\small,
    text=PTLabel,
    align=center,
  },
  PTTraceNode/.style={
    draw=black,
    fill=PTTrace!18,
    rounded corners=\PTCoreRadius,
    minimum width=7.2mm,
    minimum height=7.2mm,
    inner sep=1.4pt,
    outer sep=0pt,
    line width=\PTCoreStroke,
    font=\small,
    text=PTLabel,
    align=center,
  },
  PTDummyBond/.style={
    circle,
    fill=black,
    draw=none,
    inner sep=1.15pt,
    outer sep=0pt,
  },
}

% -----------------------------------------------------------------------------
% Index labels: math + ITensor tag + optional prime level
% -----------------------------------------------------------------------------
\makeatletter
\define@key{ptindex}{anchor}{\def\PTIndexAnchor{#1}}
\define@key{ptindex}{math}{\def\PTIndexMath{#1}}
\define@key{ptindex}{tag}{\def\PTIndexTag{#1}}
\define@key{ptindex}{prime}{\def\PTIndexPrime{#1}}
\define@key{ptindex}{xshift}{\def\PTIndexX{#1}}
\define@key{ptindex}{yshift}{\def\PTIndexY{#1}}
\makeatother

% Append a visual prime level to math (') and to a tag (ASCII apostrophe).
\newcommand{\PTIndexPrimeTokens}{%
  \ifx\PTIndexPrime\empty
  \else
    \ifnum\PTIndexPrime>0 '\fi
    \ifnum\PTIndexPrime>1 '\fi
    \ifnum\PTIndexPrime>2 '\fi
  \fi
}

\newcommand{\PTIndexLabel}[2][]{%
  \def\PTIndexAnchor{south}%
  \def\PTIndexMath{}%
  \def\PTIndexTag{}%
  \def\PTIndexPrime{}%
  \def\PTIndexX{0pt}%
  \def\PTIndexY{0pt}%
  \setkeys{ptindex}{#1}%
  \node[
    inner sep=0.9pt,
    outer sep=1.2pt,
    anchor=\PTIndexAnchor,
    align=center,
    xshift=\PTIndexX,
    yshift=\PTIndexY,
    text=PTLabel,
  ] at (#2) {%
    \def\PTIndexMathEmpty{}%
    \ifx\PTIndexMath\PTIndexMathEmpty
    \else
      {\small\(\PTIndexMath\PTIndexPrimeTokens\)}%
      \ifx\PTIndexTag\PTIndexMathEmpty
      \else \\[-0.15ex]%
      \fi
    \fi
    \ifx\PTIndexTag\PTIndexMathEmpty
    \else
      {\scriptsize\ttfamily \PTIndexTag\PTIndexPrimeTokens}%
    \fi
  };%
}

% -----------------------------------------------------------------------------
% Ket and bra boundary tensors
% -----------------------------------------------------------------------------
% Invisible rectangular node provides named anchors.  The triangle is drawn
% on top so that .west (ket) / .east (bra) is the wire attachment (base).

\newcommand{\PTKet}[3]{%
  \node[
    name=#1,
    inner sep=0pt,
    outer sep=0pt,
    minimum width=\PTKetWidth,
    minimum height=\PTKetHeight,
    font=\small,
  ] at #2 {};
  \begin{scope}[on PTforeground layer]
    \filldraw[
      draw=black,
      fill=PTHilbertKet!18,
      line width=\PTCoreStroke,
      line join=round,
    ] (#1.west |- #1.south) -- (#1.east) -- (#1.west |- #1.north) -- cycle;
    \node[font=\small, text=PTLabel, inner sep=0.4pt]
      at ($ (#1.west)!0.36!(#1.east) $) {#3};
  \end{scope}
}

\newcommand{\PTBra}[3]{%
  \node[
    name=#1,
    inner sep=0pt,
    outer sep=0pt,
    minimum width=\PTKetWidth,
    minimum height=\PTKetHeight,
    font=\small,
  ] at #2 {};
  \begin{scope}[on PTforeground layer]
    \filldraw[
      draw=black,
      fill=PTHilbertBra!18,
      line width=\PTCoreStroke,
      line join=round,
    ] (#1.east |- #1.south) -- (#1.west) -- (#1.east |- #1.north) -- cycle;
    \node[font=\small, text=PTLabel, inner sep=0.4pt]
      at ($ (#1.east)!0.36!(#1.west) $) {#3};
  \end{scope}
}

\newcommand{\PTLiouvilleKet}[3]{%
  \node[
    name=#1,
    inner sep=0pt,
    outer sep=0pt,
    minimum width=\PTKetWidth,
    minimum height=\PTKetHeight,
    font=\small,
  ] at #2 {};
  \begin{scope}[on PTforeground layer]
    \filldraw[
      draw=black,
      fill=PTLiouville!16,
      line width=\PTCoreStroke,
      line join=round,
    ] (#1.west |- #1.south) -- (#1.east) -- (#1.west |- #1.north) -- cycle;
    \node[font=\small, text=PTLabel, inner sep=0.4pt]
      at ($ (#1.west)!0.36!(#1.east) $) {#3};
  \end{scope}
}

\newcommand{\PTLiouvilleBra}[3]{%
  \node[
    name=#1,
    inner sep=0pt,
    outer sep=0pt,
    minimum width=\PTKetWidth,
    minimum height=\PTKetHeight,
    font=\small,
  ] at #2 {};
  \begin{scope}[on PTforeground layer]
    \filldraw[
      draw=black,
      fill=PTLiouville!16,
      line width=\PTCoreStroke,
      line join=round,
    ] (#1.east |- #1.south) -- (#1.west) -- (#1.east |- #1.north) -- cycle;
    \node[font=\small, text=PTLabel, inner sep=0.4pt]
      at ($ (#1.east)!0.36!(#1.west) $) {#3};
  \end{scope}
}

% -----------------------------------------------------------------------------
% Square tensor cores (thick stroke, slight corner rounding)
% -----------------------------------------------------------------------------
\makeatletter
\define@key{pttensor}{fill}{\def\PTTensorFill{#1}}
\define@key{pttensor}{draw}{\def\PTTensorDrawCol{#1}}
\define@key{pttensor}{minimum width}{\def\PTTensorMinW{#1}}
\define@key{pttensor}{minimum height}{\def\PTTensorMinH{#1}}
\define@key{pttensor}{rounded}{\def\PTTensorRnd{#1}}
\define@key{pttensor}{font}{\def\PTTensorFont{#1}}
\makeatother

\newcommand{\PTTensor}[4][]{%
  \def\PTTensorFill{PTTensor}%
  \def\PTTensorDrawCol{black}%
  \def\PTTensorMinW{\PTCoreWidth}%
  \def\PTTensorMinH{\PTCoreHeight}%
  \def\PTTensorRnd{\PTCoreRadius}%
  \def\PTTensorFont{\small}%
  \setkeys{pttensor}{#1}%
  \begin{scope}[on PTforeground layer]
    \node[
      name=#2,
      PTCore,
      fill=\PTTensorFill,
      draw=\PTTensorDrawCol,
      minimum width=\PTTensorMinW,
      minimum height=\PTTensorMinH,
      rounded corners=\PTTensorRnd,
      font=\PTTensorFont,
    ] at #3 {#4};
  \end{scope}
}

\newcommand{\PTOperator}[4][]{%
  \PTTensor[fill=PTTensor, minimum width=\PTCoreWidth, minimum height=\PTCoreWidth, #1]{#2}{#3}{#4}%
}

\newcommand{\PTSuperoperator}[4][]{%
  \PTTensor[
    fill=PTTensor,
    minimum width=\PTCoreWidth,
    minimum height=\PTCoreWidth,
    #1
  ]{#2}{#3}{#4}%
}

\newcommand{\PTProcessCore}[4][]{%
  \PTTensor[
    fill=PTProcess,
    minimum width=\PTCoreWidth,
    minimum height=\PTCoreWidth,
    #1
  ]{#2}{#3}{#4}%
}

\newcommand{\PTInstrument}[4][]{%
  \PTTensor[
    fill=PTInstrument,
    minimum width=\PTCoreWidth,
    minimum height=\PTCoreWidth,
    #1
  ]{#2}{#3}{#4}%
}

% Vertical Hilbert ket (north) and bra (south) stubs on an operator.
\newcommand{\PTOperatorLegs}[1]{%
  \draw[PTHilbertKetLeg] (#1.north) -- ++(0,\PTLegLength) coordinate (#1-ket);
  \draw[PTHilbertBraLeg] (#1.south) -- ++(0,-\PTLegLength) coordinate (#1-bra);
}

% -----------------------------------------------------------------------------
% Combiners: s \otimes s'  -->  \ell    and the inverse
% -----------------------------------------------------------------------------
\makeatletter
\define@key{ptcomb}{direction}{\def\PTCombDir{#1}}
\makeatother
\def\PTCombDirUp{up}
\def\PTCombDirDown{down}
\def\PTCombDirLeft{left}
\def\PTCombDirRight{right}

\newcommand{\PTCombiner}[4][]{%
  \def\PTCombDir{up}%
  \setkeys{ptcomb}{#1}%
  \ifx\PTCombDir\PTCombDirDown
    \def\PTCombRot{180}%
  \else\ifx\PTCombDir\PTCombDirLeft
    \def\PTCombRot{90}%
  \else\ifx\PTCombDir\PTCombDirRight
    \def\PTCombRot{270}%
  \else
    \def\PTCombRot{0}%
  \fi\fi\fi
  \begin{scope}[on PTforeground layer]
    \node[
      name=#2,
      isosceles triangle,
      isosceles triangle apex angle=60,
      shape border rotate=\PTCombRot,
      inner sep=0pt,
      outer sep=0pt,
      minimum width=\PTCombinerSize,
      draw=black,
      fill=white,
      line width=\PTCoreStroke,
      font=\scriptsize,
      text=PTLabel,
    ] at #3 {};
    \node[font=\scriptsize, text=PTLabel, inner sep=0.8pt, anchor=east]
      at ([xshift=-1.6pt]#2.west) {#4};
  \end{scope}
  \coordinate (#2-fused) at (#2.apex);
  \coordinate (#2-ket)   at (#2.left corner);
  \coordinate (#2-bra)   at (#2.right corner);
}

\newcommand{\PTCombinerInv}[4][]{%
  \PTCombiner[direction=down, #1]{#2}{#3}{#4}%
}

% -----------------------------------------------------------------------------
% Cups, caps, time arrow, callouts
% -----------------------------------------------------------------------------
% Cup (union, opens upward): used for traces and vectorisation bends.
% The two endpoints are the open wire tips; control points sit below them.
\newcommand{\PTCup}[3][]{%
  \draw[#1] (#2)
    .. controls ($ (#2) + (0,-\PTBendDepth) $)
            and ($ (#3) + (0,-\PTBendDepth) $) ..
    (#3);
}

% Trace: left and right semicircles joined by a straight bar BELOW the core
% (never behind it).  #2-loopmid is the bar midpoint for a single index tag.
% Drawn on the main layer so the bar is not covered by the core; the core
% still sits in front at the ports.
\newcommand{\PTTraceLoop}[2][PTWire]{%
  \draw[#1]
    (#2.west)
    arc[start angle=90, delta angle=180, radius=4.0mm]
    -- ([yshift=-8.0mm]#2.east)
    arc[start angle=270, delta angle=180, radius=4.0mm];
  \coordinate (#2-loopmid) at ([yshift=-8.0mm]#2.center);
}

% Transpose bends: left (ket) port over the top towards the right, now a
% primed bra (coordinate #1-i); right (bra) port under the bottom towards
% the left, now an unprimed ket (coordinate #1-j).
\newcommand{\PTTransposeBends}[1]{%
  \draw[PTHilbertBraLeg]
    (#1.west)
    arc[start angle=270, delta angle=-90, radius=3.3mm]
    arc[start angle=180, delta angle=-90, radius=3.3mm]
    coordinate (#1-itop);
  \draw[PTHilbertBraLeg primed]
    (#1-itop) -- ++(11.5mm,0) coordinate (#1-i);
  \draw[PTHilbertKetLeg]
    (#1.east)
    arc[start angle=90, delta angle=-90, radius=3.3mm]
    arc[start angle=0, delta angle=-90, radius=3.3mm]
    -- ++(-11.5mm,0) coordinate (#1-j);
}
\newcommand{\PTCap}[3][]{%
  \draw[#1] (#2)
    .. controls ($ (#2) + (0,\PTBendDepth) $)
            and ($ (#3) + (0,\PTBendDepth) $) ..
    (#3);
}

% Sideways cup used when bending a vertical bra wire into a second output.
% #2 = start (usually the operator's south / bra port).
% #3 = end   (the new upward ket-like port for the bra index).
% The bend goes to the LEFT of both points by default (column-major: |k>
% becomes the left factor).
\newcommand{\PTVecBend}[3][]{%
  \draw[#1] (#2)
    .. controls ($ (#2) + (0,-\PTBendDepth) $)
            and ($ (#3) + (0,-\PTBendDepth) $) ..
    (#3);
}

% Time arrow: ALWAYS left-pointing.  #1 = right coord, #2 = left coord.
% Optional first argument is a TikZ node-anchor for the label (default below).
\newcommand{\PTTimeArrow}[4][below=2.4pt]{%
  \draw[
    -{Stealth[length=2.2mm, width=1.55mm]},
    line width=1.05pt,
    draw=PTLabel,
    line cap=round,
  ] (#2) -- (#3)
    node[midway, #1, font=\small, text=PTLabel, inner sep=0.6pt] {#4};
}

\newcommand{\PTPanelLabel}[2]{%
  \node[
    font=\small\bfseries,
    text=PTLabel,
    inner sep=0.4pt,
    anchor=south west,
  ] at #1 {#2};
}

\newcommand{\PTCallout}[3][]{%
  \node[
    font=\scriptsize,
    text=PTLabel,
    inner sep=0.8pt,
    align=left,
    #1,
  ] at (#2) {#3};
}

\newcommand{\PTMapsTo}[2][]{%
  \node[font=\small, text=PTLabel, inner sep=1pt, #1] at #2 {\(\longmapsto\)};
}

% -----------------------------------------------------------------------------
% Instruments
% -----------------------------------------------------------------------------
\newif\ifPTInstClassical
\PTInstClassicalfalse
\def\PTInstKindBox{box}
\def\PTInstKindKet{ket}
\def\PTInstKindBra{bra}
\def\PTInstKindTrace{trace}

\makeatletter
\define@key{ptinst}{fill}{\def\PTInstFill{#1}}
\define@key{ptinst}{draw}{\def\PTInstDraw{#1}}
\define@key{ptinst}{minimum width}{\def\PTInstMinW{#1}}
\define@key{ptinst}{minimum height}{\def\PTInstMinH{#1}}
\define@key{ptinst}{kind}{\def\PTInstKind{#1}}
\define@key{ptinst}{classical}{\def\PTInstClassDir{#1}}
\define@key{ptinst}{classical label}{\def\PTInstClassLab{#1}}
\define@key{ptinst}{classical length}{\def\PTInstClassLen{#1}}
\makeatother

% Generic constructor.  Quantum legs are NOT auto-drawn: connect to the
% node anchors (north, south, east, west, or the triangle bases).
% Classical legs ARE auto-drawn when classical != none.
\newcommand{\PTInstrumentNode}[4][]{%
  \def\PTInstFill{PTInstrument}%
  \def\PTInstDraw{black}%
  \def\PTInstMinW{\PTCoreWidth}%
  \def\PTInstMinH{\PTCoreWidth}%
  \def\PTInstKind{box}%
  \def\PTInstClassDir{none}%
  \def\PTInstClassLab{}%
  \def\PTInstClassLen{6.2mm}%
  \setkeys{ptinst}{#1}%
  \def\PTInstNone{none}%
  \ifx\PTInstKind\PTInstKindKet
    \node[
      name=#2,
      inner sep=0pt,
      outer sep=0pt,
      minimum width=\PTKetWidth,
      minimum height=\PTKetHeight,
    ] at #3 {};
    \begin{scope}[on PTforeground layer]
      \filldraw[
        draw=black,
        fill=\PTInstFill,
        line width=\PTCoreStroke,
        line join=round,
      ] (#2.west |- #2.south) -- (#2.east) -- (#2.west |- #2.north) -- cycle;
      \node[font=\small, text=PTLabel, inner sep=0.4pt]
        at ($ (#2.west)!0.36!(#2.east) $) {#4};
    \end{scope}
  \else
    \ifx\PTInstKind\PTInstKindBra
      \node[
        name=#2,
        inner sep=0pt,
        outer sep=0pt,
        minimum width=\PTKetWidth,
        minimum height=\PTKetHeight,
      ] at #3 {};
      \begin{scope}[on PTforeground layer]
        \filldraw[
          draw=black,
          fill=\PTInstFill,
          line width=\PTCoreStroke,
          line join=round,
        ] (#2.east |- #2.south) -- (#2.west) -- (#2.east |- #2.north) -- cycle;
        \node[font=\small, text=PTLabel, inner sep=0.4pt]
          at ($ (#2.east)!0.36!(#2.west) $) {#4};
      \end{scope}
    \else
      \ifx\PTInstKind\PTInstKindTrace
        \begin{scope}[on PTforeground layer]
          \node[name=#2, PTTraceNode] at #3 {#4};
        \end{scope}
      \else
        \begin{scope}[on PTforeground layer]
          \node[
            name=#2,
            PTCore,
            fill=\PTInstFill,
            draw=\PTInstDraw,
            minimum width=\PTInstMinW,
            minimum height=\PTInstMinH,
          ] at #3 {#4};
        \end{scope}
      \fi
    \fi
  \fi
  \ifx\PTInstClassDir\PTInstNone
  \else
    \def\PTInstDirDown{down}%
    \def\PTInstDirUp{up}%
    \def\PTInstDirLeft{left}%
    \def\PTInstDirRight{right}%
    \ifx\PTInstClassDir\PTInstDirDown
      \draw[PTClassicalLeg] (#2.south) -- ++(0,-\PTInstClassLen)
        coordinate (#2-classical);
    \else\ifx\PTInstClassDir\PTInstDirUp
      \draw[PTClassicalLeg] (#2.north) -- ++(0,\PTInstClassLen)
        coordinate (#2-classical);
    \else\ifx\PTInstClassDir\PTInstDirLeft
      \draw[PTClassicalLeg] (#2.west) -- ++(-\PTInstClassLen,0)
        coordinate (#2-classical);
    \else
      \draw[PTClassicalLeg] (#2.east) -- ++(\PTInstClassLen,0)
        coordinate (#2-classical);
    \fi\fi\fi
    \def\PTInstClassEmpty{}%
    \ifx\PTInstClassLab\PTInstClassEmpty
    \else
      \node[font=\small, text=PTLabel, inner sep=0.8pt, anchor=north]
        at ([yshift=-1.0pt]#2-classical) {\(\PTInstClassLab\)};
    \fi
  \fi
}

\newcommand{\PTIdentity}[4][]{%
  \PTInstrumentNode[fill=PTTensor, #1]{#2}{#3}{#4}%
}

\newcommand{\PTControl}[4][]{%
  \PTInstrumentNode[#1]{#2}{#3}{#4}%
}

\newcommand{\PTUnitary}[4][]{%
  \PTInstrumentNode[#1]{#2}{#3}{#4}%
}

\newcommand{\PTPrepare}[4][]{%
  \PTInstrumentNode[kind=ket, fill=PTInstrument, #1]{#2}{#3}{#4}%
}

\newcommand{\PTMeasure}[4][]{%
  \PTInstrumentNode[
    classical=down,
    classical label={x},
    #1
  ]{#2}{#3}{#4}%
}

\newcommand{\PTProjective}[4][]{%
  \PTInstrumentNode[
    classical=down,
    classical label={x},
    #1
  ]{#2}{#3}{#4}%
}

% Measurement (earlier, right) then preparation (later, left).
% Quantum wire is broken; the classical outcome is drawn separately.
% Nodes: <name>-E, <name>-P, <name>-x.
\newcommand{\PTCausalBreak}[3][]{%
  \begin{scope}[shift={#3}]
    \PTInstrumentNode[
      classical=none,
      minimum width=8.4mm,
      minimum height=8.4mm,
      #1
    ]{#2-E}{(0,0)}{\(E^{(x)}\)}
    \PTPrepare{#2-P}{(-1.65,0)}{\(\mathcal{P}\)}
    \draw[PTClassicalLeg] (#2-E.south) -- ++(0,-5.8mm) coordinate (#2-x);
    \node[font=\small, text=PTLabel, inner sep=0.6pt, anchor=north]
      at ([yshift=-0.8pt]#2-x) {\(x\)};
  \end{scope}
}

% Hilbert-space trace: magenta cup closing ket and bra tips.
\newcommand{\PTTraceHilbert}[2]{%
  \PTCup[PTTraceCup]{#1}{#2}%
  \node[font=\small, text=PTTrace, inner sep=0.5pt, anchor=north]
    at ($ (#1)!0.5!(#2) + (0,-\PTBendDepth-0.5pt) $) {\(\mathrm{Tr}\)};
}

% Liouville-space trace: compact Tr node sitting on a fused leg.
\newcommand{\PTTraceLiouville}[3]{%
  \begin{scope}[on PTforeground layer]
    \node[name=#1, PTTraceNode] at #2 {#3};
  \end{scope}
}

% -----------------------------------------------------------------------------
% PT-MPO / joint process tensor
% -----------------------------------------------------------------------------
\newif\ifPTMPOjoint
\newif\ifPTMPOtags
\newif\ifPTMPOchi
\newif\ifPTMPOtimes
\newif\ifPTMPOarrow
\newif\ifPTMPOfinal
\PTMPOjointfalse
\PTMPOtagstrue
\PTMPOchitrue
\PTMPOtimestrue
\PTMPOarrowtrue
\PTMPOfinalfalse

\makeatletter
\define@key{ptmpo}{n}{\def\PTMPOn{#1}}
\define@key{ptmpo}{at}{\def\PTMPOat{#1}}
\define@key{ptmpo}{name}{\def\PTMPOname{#1}}
\define@key{ptmpo}{spacing}{\setlength{\PTMPOspacing}{#1}}
\define@key{ptmpo}{joint}{\csname PTMPOjoint#1\endcsname}
\define@key{ptmpo}{tags}{\csname PTMPOtags#1\endcsname}
\define@key{ptmpo}{chi}{\csname PTMPOchi#1\endcsname}
\define@key{ptmpo}{times}{\csname PTMPOtimes#1\endcsname}
\define@key{ptmpo}{time arrow}{\csname PTMPOarrow#1\endcsname}
\define@key{ptmpo}{final}{\csname PTMPOfinal#1\endcsname}
\define@key{ptmpo}{label}{\def\PTMPOlabel{#1}}
\makeatother

% Draw n time cores W^{[0]} (right, t_0) ... W^{[n-1]} (left).
% Node names:  <name>-W-<k>
% IO tips:     <name>-out-<k>  (process output, RIGHT of pair)
%              <name>-in-<k>   (process input,  LEFT  of pair)
% Time arrow:  <name>-t0  (right)  -->  <name>-tn (left)
\newcommand{\PTDrawProcessMPO}[1][]{%
  \PTMPOjointfalse
  \PTMPOtagstrue
  \PTMPOchitrue
  \PTMPOtimestrue
  \PTMPOarrowtrue
  \PTMPOfinalfalse
  \setlength{\PTMPOspacing}{\PTTimeSpacing}%
  \def\PTMPOn{3}%
  \def\PTMPOat{0,0}%
  \def\PTMPOname{pt}%
  \def\PTMPOlabel{}%
  \setkeys{ptmpo}{#1}%
  \coordinate (\PTMPOname-origin) at (\PTMPOat);
  \edef\PTMPOlast{\PTMPOname-W-\the\numexpr\PTMPOn-1\relax}%
  % --- cores, right (k=0) to left ------------------------------------------
  \foreach \k in {0,...,\the\numexpr\PTMPOn-1\relax} {%
    \path (\PTMPOname-origin)
      ++(-\k*\PTMPOspacing,0)
      coordinate (\PTMPOname-c-\k);
    \PTProcessCore{\PTMPOname-W-\k}{(\PTMPOname-c-\k)}{\(W^{[\k]}\)}%
    % Local ports: output on the right, input on the left (local time).
    \coordinate (\PTMPOname-outsrc-\k) at
      ([xshift=\PTIOSep]\PTMPOname-W-\k.south);
    \coordinate (\PTMPOname-insrc-\k) at
      ([xshift=-\PTIOSep]\PTMPOname-W-\k.south);
    \draw[PTLiouvilleLeg]
      (\PTMPOname-outsrc-\k) -- ++(0,-\PTLegLength)
      coordinate (\PTMPOname-out-\k);
    \draw[PTLiouvilleLeg]
      (\PTMPOname-insrc-\k) -- ++(0,-\PTLegLength)
      coordinate (\PTMPOname-in-\k);
  }%
  % --- joint body (no exposed link indices) or MPO links -------------------
  \ifPTMPOjoint
    \node[
      fit=(\PTMPOname-W-0) (\PTMPOlast),
      inner xsep=3.6pt,
      inner ysep=4.6pt,
      rounded corners=5.5pt,
      fill=PTProcess!55,
      draw=black,
      line width=0.7pt,
    ] (\PTMPOname-body) {};
    \draw[
      draw=PTProcess!70!black,
      line width=3.4pt,
      line cap=round,
    ] (\PTMPOlast.center) -- (\PTMPOname-W-0.center);
  \else
    \ifnum\PTMPOn>1
      \foreach \k [
        evaluate=\k as \knext using int(\k+1)
      ] in {0,...,\the\numexpr\PTMPOn-2\relax} {%
        \draw[PTLinkLeg]
          (\PTMPOname-W-\k.west) -- (\PTMPOname-W-\knext.east)
          coordinate[midway] (\PTMPOname-link-\k);
        \ifPTMPOchi
          \node[font=\small, text=PTLabel, inner sep=0.5pt, above=1.4pt]
            at (\PTMPOname-link-\k) {\(\chi_{\k}\)};
        \fi
        \ifPTMPOtags
          \PTIndexLabel[
            anchor=south,
            tag={PTLink,t-0\k},
            yshift=11.0pt,
          ]{\PTMPOname-link-\k}%
        \fi
      }%
    \fi
    % Dimension-1 boundary stubs.
    \draw[PTLinkLeg] (\PTMPOname-W-0.east) -- ++(4.2mm,0)
      coordinate (\PTMPOname-bnd-R);
    \begin{scope}[on PTforeground layer]
      \node[PTDummyBond] (\PTMPOname-dummy-R) at (\PTMPOname-bnd-R) {};
    \end{scope}
    \draw[PTLinkLeg]
      (\PTMPOlast.west) -- ++(-4.2mm,0)
      coordinate (\PTMPOname-bnd-L);
    \begin{scope}[on PTforeground layer]
      \node[PTDummyBond] (\PTMPOname-dummy-L) at (\PTMPOname-bnd-L) {};
    \end{scope}
  \fi
  % --- optional final (t_n) physical output on the far left ----------------
  \ifPTMPOfinal
    \coordinate (\PTMPOname-final-src) at ([yshift=1.6mm]\PTMPOlast.west);
    \draw[PTLiouvilleLeg] (\PTMPOname-final-src) -- ++(-7.5mm,0)
      coordinate (\PTMPOname-final);
  \fi
  % --- time labels and left-pointing arrow ---------------------------------
  \coordinate (\PTMPOname-t0) at ([yshift=-15.2mm, xshift=5.0mm]\PTMPOname-W-0.south);
  \coordinate (\PTMPOname-tn) at ([yshift=-15.2mm, xshift=-5.0mm]\PTMPOlast.south);
  \ifPTMPOtimes
    \node[font=\small, text=PTLabel, inner sep=0.4pt, anchor=west]
      at (\PTMPOname-t0) {\(t_{0}\)};
    \node[font=\small, text=PTLabel, inner sep=0.4pt, anchor=east]
      at (\PTMPOname-tn) {\(t_{n}\)};
  \fi
  \ifPTMPOarrow
    \PTTimeArrow[below=1.2pt]{\PTMPOname-t0}{\PTMPOname-tn}{\(t\)}%
  \fi
  \ifPTMPOtags
    \PTIndexLabel[
      anchor=west,
      math={\ell_{\mathrm{out}}},
      tag={Output,t-00},
      xshift=3.0pt,
    ]{\PTMPOname-out-0}%
    \PTIndexLabel[
      anchor=east,
      math={\ell_{\mathrm{in}}},
      tag={Input,t-00},
      xshift=-3.0pt,
      yshift=-4.2pt,
    ]{\PTMPOname-in-0}%
  \fi
  \def\PTMPOlabelEmpty{}%
  \ifx\PTMPOlabel\PTMPOlabelEmpty
  \else
    \node[font=\small, text=PTLabel, inner sep=1pt, above=3.6pt]
      at ($ (\PTMPOname-W-0.north)!0.5!(\PTMPOlast.north) $)
      {\PTMPOlabel};
  \fi
}

\newcommand{\PTDrawJointProcess}[1][]{%
  \PTDrawProcessMPO[joint=true,chi=false,tags=false,#1]%
}

% -----------------------------------------------------------------------------
% Example of a future instrument (commented; copy and adapt).
% -----------------------------------------------------------------------------
% \newcommand{\PTReset}[4][]{%
%   \PTInstrumentNode[fill=PTInstrument, classical=none, #1]{#2}{#3}{#4}%
% }
% % Tester / ancilla interaction (two-site box):
% \newcommand{\PTTester}[4][]{%
%   \PTInstrumentNode[
%     fill=PTInstrument,
%     minimum width=1.55\PTCoreWidth,
%     #1
%   ]{#2}{#3}{#4}%
% }
% % Feedback-conditioned operation: reuse classical=down and wire it into
% % a later \PTInstrumentNode.
%
% This file is a library.  Compile pt_tikz_demo.tex for a catalogue of
% every primitive.

% Figure fragments then:
%   \input{figures/tikz/<name>.tikz}
%
% =============================================================================

\ifdefined\PTTikZLoaded
  \expandafter\endinput
\fi
\def\PTTikZLoaded{1}

\RequirePackage{tikz}
\RequirePackage{xcolor}
\RequirePackage{keyval}
\usetikzlibrary{
  arrows.meta,
  backgrounds,
  calc,
  decorations.markings,
  fit,
  matrix,
  positioning,
  shapes.geometric
}

% -----------------------------------------------------------------------------
% Colours (Okabe--Ito / colour-blind friendly; intelligible in grayscale)
% -----------------------------------------------------------------------------
\definecolor{PTHilbertKet}{HTML}{0072B2}
\definecolor{PTHilbertBra}{HTML}{D55E00}
\definecolor{PTLiouville}{HTML}{7B3294}
\definecolor{PTLink}{HTML}{009E73}
\definecolor{PTTensor}{HTML}{E8E8E8}
\definecolor{PTProcess}{HTML}{DCEAF7}
\definecolor{PTInstrument}{HTML}{F6E8A6}
\definecolor{PTTrace}{HTML}{CC79A7}
\definecolor{PTMeasurement}{HTML}{E69F00}
\definecolor{PTLabel}{HTML}{333333}
\definecolor{PTBond}{HTML}{6A6A6A}

% -----------------------------------------------------------------------------
% Lengths (the single place to retune the visual language)
% -----------------------------------------------------------------------------
\newlength{\PTCoreWidth}      \setlength{\PTCoreWidth}{9.6mm}
\newlength{\PTCoreHeight}     \setlength{\PTCoreHeight}{9.6mm}
\newlength{\PTTimeSpacing}    \setlength{\PTTimeSpacing}{20.0mm}
\newlength{\PTLegLength}      \setlength{\PTLegLength}{7.0mm}
\newlength{\PTIOSep}          \setlength{\PTIOSep}{2.4mm}
% Equilateral triangle: vertical base = side, horizontal altitude = (sqrt(3)/2)*side.
\newlength{\PTKetHeight}      \setlength{\PTKetHeight}{7.8mm}
\newlength{\PTKetWidth}       \setlength{\PTKetWidth}{6.75mm}
\newlength{\PTBendDepth}      \setlength{\PTBendDepth}{5.5mm}
\newlength{\PTCombinerSize}   \setlength{\PTCombinerSize}{7.2mm}
\newlength{\PTHilbertLW}      \setlength{\PTHilbertLW}{1.35pt}
\newlength{\PTLiouvilleLW}    \setlength{\PTLiouvilleLW}{0.70pt}
\newlength{\PTLiouvilleGap}   \setlength{\PTLiouvilleGap}{1.55pt}
\newlength{\PTLinkLW}         \setlength{\PTLinkLW}{2.70pt}
\newlength{\PTBondLW}         \setlength{\PTBondLW}{1.35pt}
\newlength{\PTClassicalLW}    \setlength{\PTClassicalLW}{1.15pt}
\newlength{\PTCoreStroke}     \setlength{\PTCoreStroke}{1.45pt}
\newlength{\PTCoreRadius}     \setlength{\PTCoreRadius}{1.3pt}
\newlength{\PTBridgeHalo}     \setlength{\PTBridgeHalo}{4.2pt}

% Scratch length for the MPO drawer (not a public tuning knob).
\newlength{\PTMPOspacing}

% -----------------------------------------------------------------------------
% Layers: wires behind, labels in the middle, filled shapes in front.
% -----------------------------------------------------------------------------
\pgfdeclarelayer{PTforeground}
\pgfsetlayers{background,main,PTforeground}
\tikzset{
  on PTforeground layer/.style={
    execute at begin scope={\pgfonlayer{PTforeground}},
    execute at end scope={\endpgfonlayer},
  },
}

% -----------------------------------------------------------------------------
% Line and node styles
% -----------------------------------------------------------------------------
\tikzset{
  PTPicture/.style={
    baseline,
    line join=round,
    >={Stealth[length=2.0mm, width=1.45mm]},
  },
  % --- wires (black; structure, not colour, distinguishes spaces) ----------
  PTWire/.style={
    draw=black,
    line width=\PTHilbertLW,
    line cap=butt,
    line join=round,
    double=none,
  },
  PTHilbertKetLeg/.style={
    PTWire,
    on background layer,
  },
  PTHilbertBraLeg/.style={
    draw=black,
    line width=\PTHilbertLW,
    line cap=butt,
    line join=round,
    double=none,
    on background layer,
  },
  % Optional prime tick at the midpoint of a bra wire.
  PTHilbertBraLeg primed/.style={
    PTHilbertBraLeg,
    postaction={
      decorate,
      decoration={
        markings,
        mark=at position 0.55 with {
          \node[font=\scriptsize, text=PTLabel, inner sep=0.4pt,
                xshift=3.2pt, yshift=0.4pt] {\(\prime\)};
        }
      }
    },
  },
  PTLiouvilleLeg/.style={
    draw=black,
    line width=\PTLiouvilleLW,
    double=white,
    double distance=\PTLiouvilleGap,
    line cap=butt,
    line join=round,
    on background layer,
  },
  PTBondLeg/.style={
    draw=black,
    line width=\PTBondLW,
    line cap=butt,
    line join=round,
    double=none,
    on background layer,
  },
  PTLinkLeg/.style={
    draw=black,
    line width=\PTLinkLW,
    line cap=butt,
    line join=round,
    double=none,
    on background layer,
  },
  PTClassicalLeg/.style={
    draw=black,
    line width=\PTClassicalLW,
    dash pattern=on 2.4pt off 1.6pt,
    line cap=butt,
    line join=round,
    double=none,
    on background layer,
  },
  PTTraceCup/.style={
    draw=PTTrace,
    line width=\PTHilbertLW,
    line cap=round,
    line join=round,
    double=none,
    on background layer,
  },
  % White halo so an over-crossing does not look like a contraction.
  PTBridge/.style={
    preaction={
      draw=white,
      line width=\PTBridgeHalo,
      double=none,
      line cap=butt,
    },
  },
  % --- cores ---------------------------------------------------------------
  PTCore/.style={
    draw=black,
    fill=PTTensor,
    rounded corners=\PTCoreRadius,
    minimum width=\PTCoreWidth,
    minimum height=\PTCoreHeight,
    inner sep=2.2pt,
    outer sep=0pt,
    line width=\PTCoreStroke,
    font=\small,
    text=PTLabel,
    align=center,
  },
  PTTraceNode/.style={
    draw=black,
    fill=PTTrace!18,
    rounded corners=\PTCoreRadius,
    minimum width=7.2mm,
    minimum height=7.2mm,
    inner sep=1.4pt,
    outer sep=0pt,
    line width=\PTCoreStroke,
    font=\small,
    text=PTLabel,
    align=center,
  },
  PTDummyBond/.style={
    circle,
    fill=black,
    draw=none,
    inner sep=1.15pt,
    outer sep=0pt,
  },
}

% -----------------------------------------------------------------------------
% Index labels: math + ITensor tag + optional prime level
% -----------------------------------------------------------------------------
\makeatletter
\define@key{ptindex}{anchor}{\def\PTIndexAnchor{#1}}
\define@key{ptindex}{math}{\def\PTIndexMath{#1}}
\define@key{ptindex}{tag}{\def\PTIndexTag{#1}}
\define@key{ptindex}{prime}{\def\PTIndexPrime{#1}}
\define@key{ptindex}{xshift}{\def\PTIndexX{#1}}
\define@key{ptindex}{yshift}{\def\PTIndexY{#1}}
\makeatother

% Append a visual prime level to math (') and to a tag (ASCII apostrophe).
\newcommand{\PTIndexPrimeTokens}{%
  \ifx\PTIndexPrime\empty
  \else
    \ifnum\PTIndexPrime>0 '\fi
    \ifnum\PTIndexPrime>1 '\fi
    \ifnum\PTIndexPrime>2 '\fi
  \fi
}

\newcommand{\PTIndexLabel}[2][]{%
  \def\PTIndexAnchor{south}%
  \def\PTIndexMath{}%
  \def\PTIndexTag{}%
  \def\PTIndexPrime{}%
  \def\PTIndexX{0pt}%
  \def\PTIndexY{0pt}%
  \setkeys{ptindex}{#1}%
  \node[
    inner sep=0.9pt,
    outer sep=1.2pt,
    anchor=\PTIndexAnchor,
    align=center,
    xshift=\PTIndexX,
    yshift=\PTIndexY,
    text=PTLabel,
  ] at (#2) {%
    \def\PTIndexMathEmpty{}%
    \ifx\PTIndexMath\PTIndexMathEmpty
    \else
      {\small\(\PTIndexMath\PTIndexPrimeTokens\)}%
      \ifx\PTIndexTag\PTIndexMathEmpty
      \else \\[-0.15ex]%
      \fi
    \fi
    \ifx\PTIndexTag\PTIndexMathEmpty
    \else
      {\scriptsize\ttfamily \PTIndexTag\PTIndexPrimeTokens}%
    \fi
  };%
}

% -----------------------------------------------------------------------------
% Ket and bra boundary tensors
% -----------------------------------------------------------------------------
% Invisible rectangular node provides named anchors.  The triangle is drawn
% on top so that .west (ket) / .east (bra) is the wire attachment (base).

\newcommand{\PTKet}[3]{%
  \node[
    name=#1,
    inner sep=0pt,
    outer sep=0pt,
    minimum width=\PTKetWidth,
    minimum height=\PTKetHeight,
    font=\small,
  ] at #2 {};
  \begin{scope}[on PTforeground layer]
    \filldraw[
      draw=black,
      fill=PTHilbertKet!18,
      line width=\PTCoreStroke,
      line join=round,
    ] (#1.west |- #1.south) -- (#1.east) -- (#1.west |- #1.north) -- cycle;
    \node[font=\small, text=PTLabel, inner sep=0.4pt]
      at ($ (#1.west)!0.36!(#1.east) $) {#3};
  \end{scope}
}

\newcommand{\PTBra}[3]{%
  \node[
    name=#1,
    inner sep=0pt,
    outer sep=0pt,
    minimum width=\PTKetWidth,
    minimum height=\PTKetHeight,
    font=\small,
  ] at #2 {};
  \begin{scope}[on PTforeground layer]
    \filldraw[
      draw=black,
      fill=PTHilbertBra!18,
      line width=\PTCoreStroke,
      line join=round,
    ] (#1.east |- #1.south) -- (#1.west) -- (#1.east |- #1.north) -- cycle;
    \node[font=\small, text=PTLabel, inner sep=0.4pt]
      at ($ (#1.east)!0.36!(#1.west) $) {#3};
  \end{scope}
}

\newcommand{\PTLiouvilleKet}[3]{%
  \node[
    name=#1,
    inner sep=0pt,
    outer sep=0pt,
    minimum width=\PTKetWidth,
    minimum height=\PTKetHeight,
    font=\small,
  ] at #2 {};
  \begin{scope}[on PTforeground layer]
    \filldraw[
      draw=black,
      fill=PTLiouville!16,
      line width=\PTCoreStroke,
      line join=round,
    ] (#1.west |- #1.south) -- (#1.east) -- (#1.west |- #1.north) -- cycle;
    \node[font=\small, text=PTLabel, inner sep=0.4pt]
      at ($ (#1.west)!0.36!(#1.east) $) {#3};
  \end{scope}
}

\newcommand{\PTLiouvilleBra}[3]{%
  \node[
    name=#1,
    inner sep=0pt,
    outer sep=0pt,
    minimum width=\PTKetWidth,
    minimum height=\PTKetHeight,
    font=\small,
  ] at #2 {};
  \begin{scope}[on PTforeground layer]
    \filldraw[
      draw=black,
      fill=PTLiouville!16,
      line width=\PTCoreStroke,
      line join=round,
    ] (#1.east |- #1.south) -- (#1.west) -- (#1.east |- #1.north) -- cycle;
    \node[font=\small, text=PTLabel, inner sep=0.4pt]
      at ($ (#1.east)!0.36!(#1.west) $) {#3};
  \end{scope}
}

% -----------------------------------------------------------------------------
% Square tensor cores (thick stroke, slight corner rounding)
% -----------------------------------------------------------------------------
\makeatletter
\define@key{pttensor}{fill}{\def\PTTensorFill{#1}}
\define@key{pttensor}{draw}{\def\PTTensorDrawCol{#1}}
\define@key{pttensor}{minimum width}{\def\PTTensorMinW{#1}}
\define@key{pttensor}{minimum height}{\def\PTTensorMinH{#1}}
\define@key{pttensor}{rounded}{\def\PTTensorRnd{#1}}
\define@key{pttensor}{font}{\def\PTTensorFont{#1}}
\makeatother

\newcommand{\PTTensor}[4][]{%
  \def\PTTensorFill{PTTensor}%
  \def\PTTensorDrawCol{black}%
  \def\PTTensorMinW{\PTCoreWidth}%
  \def\PTTensorMinH{\PTCoreHeight}%
  \def\PTTensorRnd{\PTCoreRadius}%
  \def\PTTensorFont{\small}%
  \setkeys{pttensor}{#1}%
  \begin{scope}[on PTforeground layer]
    \node[
      name=#2,
      PTCore,
      fill=\PTTensorFill,
      draw=\PTTensorDrawCol,
      minimum width=\PTTensorMinW,
      minimum height=\PTTensorMinH,
      rounded corners=\PTTensorRnd,
      font=\PTTensorFont,
    ] at #3 {#4};
  \end{scope}
}

\newcommand{\PTOperator}[4][]{%
  \PTTensor[fill=PTTensor, minimum width=\PTCoreWidth, minimum height=\PTCoreWidth, #1]{#2}{#3}{#4}%
}

\newcommand{\PTSuperoperator}[4][]{%
  \PTTensor[
    fill=PTTensor,
    minimum width=\PTCoreWidth,
    minimum height=\PTCoreWidth,
    #1
  ]{#2}{#3}{#4}%
}

\newcommand{\PTProcessCore}[4][]{%
  \PTTensor[
    fill=PTProcess,
    minimum width=\PTCoreWidth,
    minimum height=\PTCoreWidth,
    #1
  ]{#2}{#3}{#4}%
}

\newcommand{\PTInstrument}[4][]{%
  \PTTensor[
    fill=PTInstrument,
    minimum width=\PTCoreWidth,
    minimum height=\PTCoreWidth,
    #1
  ]{#2}{#3}{#4}%
}

% Vertical Hilbert ket (north) and bra (south) stubs on an operator.
\newcommand{\PTOperatorLegs}[1]{%
  \draw[PTHilbertKetLeg] (#1.north) -- ++(0,\PTLegLength) coordinate (#1-ket);
  \draw[PTHilbertBraLeg] (#1.south) -- ++(0,-\PTLegLength) coordinate (#1-bra);
}

% -----------------------------------------------------------------------------
% Combiners: s \otimes s'  -->  \ell    and the inverse
% -----------------------------------------------------------------------------
\makeatletter
\define@key{ptcomb}{direction}{\def\PTCombDir{#1}}
\makeatother
\def\PTCombDirUp{up}
\def\PTCombDirDown{down}
\def\PTCombDirLeft{left}
\def\PTCombDirRight{right}

\newcommand{\PTCombiner}[4][]{%
  \def\PTCombDir{up}%
  \setkeys{ptcomb}{#1}%
  \ifx\PTCombDir\PTCombDirDown
    \def\PTCombRot{180}%
  \else\ifx\PTCombDir\PTCombDirLeft
    \def\PTCombRot{90}%
  \else\ifx\PTCombDir\PTCombDirRight
    \def\PTCombRot{270}%
  \else
    \def\PTCombRot{0}%
  \fi\fi\fi
  \begin{scope}[on PTforeground layer]
    \node[
      name=#2,
      isosceles triangle,
      isosceles triangle apex angle=60,
      shape border rotate=\PTCombRot,
      inner sep=0pt,
      outer sep=0pt,
      minimum width=\PTCombinerSize,
      draw=black,
      fill=white,
      line width=\PTCoreStroke,
      font=\scriptsize,
      text=PTLabel,
    ] at #3 {};
    \node[font=\scriptsize, text=PTLabel, inner sep=0.8pt, anchor=east]
      at ([xshift=-1.6pt]#2.west) {#4};
  \end{scope}
  \coordinate (#2-fused) at (#2.apex);
  \coordinate (#2-ket)   at (#2.left corner);
  \coordinate (#2-bra)   at (#2.right corner);
}

\newcommand{\PTCombinerInv}[4][]{%
  \PTCombiner[direction=down, #1]{#2}{#3}{#4}%
}

% -----------------------------------------------------------------------------
% Cups, caps, time arrow, callouts
% -----------------------------------------------------------------------------
% Cup (union, opens upward): used for traces and vectorisation bends.
% The two endpoints are the open wire tips; control points sit below them.
\newcommand{\PTCup}[3][]{%
  \draw[#1] (#2)
    .. controls ($ (#2) + (0,-\PTBendDepth) $)
            and ($ (#3) + (0,-\PTBendDepth) $) ..
    (#3);
}

% Trace: left and right semicircles joined by a straight bar BELOW the core
% (never behind it).  #2-loopmid is the bar midpoint for a single index tag.
% Drawn on the main layer so the bar is not covered by the core; the core
% still sits in front at the ports.
\newcommand{\PTTraceLoop}[2][PTWire]{%
  \draw[#1]
    (#2.west)
    arc[start angle=90, delta angle=180, radius=4.0mm]
    -- ([yshift=-8.0mm]#2.east)
    arc[start angle=270, delta angle=180, radius=4.0mm];
  \coordinate (#2-loopmid) at ([yshift=-8.0mm]#2.center);
}

% Transpose bends: left (ket) port over the top towards the right, now a
% primed bra (coordinate #1-i); right (bra) port under the bottom towards
% the left, now an unprimed ket (coordinate #1-j).
\newcommand{\PTTransposeBends}[1]{%
  \draw[PTHilbertBraLeg]
    (#1.west)
    arc[start angle=270, delta angle=-90, radius=3.3mm]
    arc[start angle=180, delta angle=-90, radius=3.3mm]
    coordinate (#1-itop);
  \draw[PTHilbertBraLeg primed]
    (#1-itop) -- ++(11.5mm,0) coordinate (#1-i);
  \draw[PTHilbertKetLeg]
    (#1.east)
    arc[start angle=90, delta angle=-90, radius=3.3mm]
    arc[start angle=0, delta angle=-90, radius=3.3mm]
    -- ++(-11.5mm,0) coordinate (#1-j);
}
\newcommand{\PTCap}[3][]{%
  \draw[#1] (#2)
    .. controls ($ (#2) + (0,\PTBendDepth) $)
            and ($ (#3) + (0,\PTBendDepth) $) ..
    (#3);
}

% Sideways cup used when bending a vertical bra wire into a second output.
% #2 = start (usually the operator's south / bra port).
% #3 = end   (the new upward ket-like port for the bra index).
% The bend goes to the LEFT of both points by default (column-major: |k>
% becomes the left factor).
\newcommand{\PTVecBend}[3][]{%
  \draw[#1] (#2)
    .. controls ($ (#2) + (0,-\PTBendDepth) $)
            and ($ (#3) + (0,-\PTBendDepth) $) ..
    (#3);
}

% Time arrow: ALWAYS left-pointing.  #1 = right coord, #2 = left coord.
% Optional first argument is a TikZ node-anchor for the label (default below).
\newcommand{\PTTimeArrow}[4][below=2.4pt]{%
  \draw[
    -{Stealth[length=2.2mm, width=1.55mm]},
    line width=1.05pt,
    draw=PTLabel,
    line cap=round,
  ] (#2) -- (#3)
    node[midway, #1, font=\small, text=PTLabel, inner sep=0.6pt] {#4};
}

\newcommand{\PTPanelLabel}[2]{%
  \node[
    font=\small\bfseries,
    text=PTLabel,
    inner sep=0.4pt,
    anchor=south west,
  ] at #1 {#2};
}

\newcommand{\PTCallout}[3][]{%
  \node[
    font=\scriptsize,
    text=PTLabel,
    inner sep=0.8pt,
    align=left,
    #1,
  ] at (#2) {#3};
}

\newcommand{\PTMapsTo}[2][]{%
  \node[font=\small, text=PTLabel, inner sep=1pt, #1] at #2 {\(\longmapsto\)};
}

% -----------------------------------------------------------------------------
% Instruments
% -----------------------------------------------------------------------------
\newif\ifPTInstClassical
\PTInstClassicalfalse
\def\PTInstKindBox{box}
\def\PTInstKindKet{ket}
\def\PTInstKindBra{bra}
\def\PTInstKindTrace{trace}

\makeatletter
\define@key{ptinst}{fill}{\def\PTInstFill{#1}}
\define@key{ptinst}{draw}{\def\PTInstDraw{#1}}
\define@key{ptinst}{minimum width}{\def\PTInstMinW{#1}}
\define@key{ptinst}{minimum height}{\def\PTInstMinH{#1}}
\define@key{ptinst}{kind}{\def\PTInstKind{#1}}
\define@key{ptinst}{classical}{\def\PTInstClassDir{#1}}
\define@key{ptinst}{classical label}{\def\PTInstClassLab{#1}}
\define@key{ptinst}{classical length}{\def\PTInstClassLen{#1}}
\makeatother

% Generic constructor.  Quantum legs are NOT auto-drawn: connect to the
% node anchors (north, south, east, west, or the triangle bases).
% Classical legs ARE auto-drawn when classical != none.
\newcommand{\PTInstrumentNode}[4][]{%
  \def\PTInstFill{PTInstrument}%
  \def\PTInstDraw{black}%
  \def\PTInstMinW{\PTCoreWidth}%
  \def\PTInstMinH{\PTCoreWidth}%
  \def\PTInstKind{box}%
  \def\PTInstClassDir{none}%
  \def\PTInstClassLab{}%
  \def\PTInstClassLen{6.2mm}%
  \setkeys{ptinst}{#1}%
  \def\PTInstNone{none}%
  \ifx\PTInstKind\PTInstKindKet
    \node[
      name=#2,
      inner sep=0pt,
      outer sep=0pt,
      minimum width=\PTKetWidth,
      minimum height=\PTKetHeight,
    ] at #3 {};
    \begin{scope}[on PTforeground layer]
      \filldraw[
        draw=black,
        fill=\PTInstFill,
        line width=\PTCoreStroke,
        line join=round,
      ] (#2.west |- #2.south) -- (#2.east) -- (#2.west |- #2.north) -- cycle;
      \node[font=\small, text=PTLabel, inner sep=0.4pt]
        at ($ (#2.west)!0.36!(#2.east) $) {#4};
    \end{scope}
  \else
    \ifx\PTInstKind\PTInstKindBra
      \node[
        name=#2,
        inner sep=0pt,
        outer sep=0pt,
        minimum width=\PTKetWidth,
        minimum height=\PTKetHeight,
      ] at #3 {};
      \begin{scope}[on PTforeground layer]
        \filldraw[
          draw=black,
          fill=\PTInstFill,
          line width=\PTCoreStroke,
          line join=round,
        ] (#2.east |- #2.south) -- (#2.west) -- (#2.east |- #2.north) -- cycle;
        \node[font=\small, text=PTLabel, inner sep=0.4pt]
          at ($ (#2.east)!0.36!(#2.west) $) {#4};
      \end{scope}
    \else
      \ifx\PTInstKind\PTInstKindTrace
        \begin{scope}[on PTforeground layer]
          \node[name=#2, PTTraceNode] at #3 {#4};
        \end{scope}
      \else
        \begin{scope}[on PTforeground layer]
          \node[
            name=#2,
            PTCore,
            fill=\PTInstFill,
            draw=\PTInstDraw,
            minimum width=\PTInstMinW,
            minimum height=\PTInstMinH,
          ] at #3 {#4};
        \end{scope}
      \fi
    \fi
  \fi
  \ifx\PTInstClassDir\PTInstNone
  \else
    \def\PTInstDirDown{down}%
    \def\PTInstDirUp{up}%
    \def\PTInstDirLeft{left}%
    \def\PTInstDirRight{right}%
    \ifx\PTInstClassDir\PTInstDirDown
      \draw[PTClassicalLeg] (#2.south) -- ++(0,-\PTInstClassLen)
        coordinate (#2-classical);
    \else\ifx\PTInstClassDir\PTInstDirUp
      \draw[PTClassicalLeg] (#2.north) -- ++(0,\PTInstClassLen)
        coordinate (#2-classical);
    \else\ifx\PTInstClassDir\PTInstDirLeft
      \draw[PTClassicalLeg] (#2.west) -- ++(-\PTInstClassLen,0)
        coordinate (#2-classical);
    \else
      \draw[PTClassicalLeg] (#2.east) -- ++(\PTInstClassLen,0)
        coordinate (#2-classical);
    \fi\fi\fi
    \def\PTInstClassEmpty{}%
    \ifx\PTInstClassLab\PTInstClassEmpty
    \else
      \node[font=\small, text=PTLabel, inner sep=0.8pt, anchor=north]
        at ([yshift=-1.0pt]#2-classical) {\(\PTInstClassLab\)};
    \fi
  \fi
}

\newcommand{\PTIdentity}[4][]{%
  \PTInstrumentNode[fill=PTTensor, #1]{#2}{#3}{#4}%
}

\newcommand{\PTControl}[4][]{%
  \PTInstrumentNode[#1]{#2}{#3}{#4}%
}

\newcommand{\PTUnitary}[4][]{%
  \PTInstrumentNode[#1]{#2}{#3}{#4}%
}

\newcommand{\PTPrepare}[4][]{%
  \PTInstrumentNode[kind=ket, fill=PTInstrument, #1]{#2}{#3}{#4}%
}

\newcommand{\PTMeasure}[4][]{%
  \PTInstrumentNode[
    classical=down,
    classical label={x},
    #1
  ]{#2}{#3}{#4}%
}

\newcommand{\PTProjective}[4][]{%
  \PTInstrumentNode[
    classical=down,
    classical label={x},
    #1
  ]{#2}{#3}{#4}%
}

% Measurement (earlier, right) then preparation (later, left).
% Quantum wire is broken; the classical outcome is drawn separately.
% Nodes: <name>-E, <name>-P, <name>-x.
\newcommand{\PTCausalBreak}[3][]{%
  \begin{scope}[shift={#3}]
    \PTInstrumentNode[
      classical=none,
      minimum width=8.4mm,
      minimum height=8.4mm,
      #1
    ]{#2-E}{(0,0)}{\(E^{(x)}\)}
    \PTPrepare{#2-P}{(-1.65,0)}{\(\mathcal{P}\)}
    \draw[PTClassicalLeg] (#2-E.south) -- ++(0,-5.8mm) coordinate (#2-x);
    \node[font=\small, text=PTLabel, inner sep=0.6pt, anchor=north]
      at ([yshift=-0.8pt]#2-x) {\(x\)};
  \end{scope}
}

% Hilbert-space trace: magenta cup closing ket and bra tips.
\newcommand{\PTTraceHilbert}[2]{%
  \PTCup[PTTraceCup]{#1}{#2}%
  \node[font=\small, text=PTTrace, inner sep=0.5pt, anchor=north]
    at ($ (#1)!0.5!(#2) + (0,-\PTBendDepth-0.5pt) $) {\(\mathrm{Tr}\)};
}

% Liouville-space trace: compact Tr node sitting on a fused leg.
\newcommand{\PTTraceLiouville}[3]{%
  \begin{scope}[on PTforeground layer]
    \node[name=#1, PTTraceNode] at #2 {#3};
  \end{scope}
}

% -----------------------------------------------------------------------------
% PT-MPO / joint process tensor
% -----------------------------------------------------------------------------
\newif\ifPTMPOjoint
\newif\ifPTMPOtags
\newif\ifPTMPOchi
\newif\ifPTMPOtimes
\newif\ifPTMPOarrow
\newif\ifPTMPOfinal
\PTMPOjointfalse
\PTMPOtagstrue
\PTMPOchitrue
\PTMPOtimestrue
\PTMPOarrowtrue
\PTMPOfinalfalse

\makeatletter
\define@key{ptmpo}{n}{\def\PTMPOn{#1}}
\define@key{ptmpo}{at}{\def\PTMPOat{#1}}
\define@key{ptmpo}{name}{\def\PTMPOname{#1}}
\define@key{ptmpo}{spacing}{\setlength{\PTMPOspacing}{#1}}
\define@key{ptmpo}{joint}{\csname PTMPOjoint#1\endcsname}
\define@key{ptmpo}{tags}{\csname PTMPOtags#1\endcsname}
\define@key{ptmpo}{chi}{\csname PTMPOchi#1\endcsname}
\define@key{ptmpo}{times}{\csname PTMPOtimes#1\endcsname}
\define@key{ptmpo}{time arrow}{\csname PTMPOarrow#1\endcsname}
\define@key{ptmpo}{final}{\csname PTMPOfinal#1\endcsname}
\define@key{ptmpo}{label}{\def\PTMPOlabel{#1}}
\makeatother

% Draw n time cores W^{[0]} (right, t_0) ... W^{[n-1]} (left).
% Node names:  <name>-W-<k>
% IO tips:     <name>-out-<k>  (process output, RIGHT of pair)
%              <name>-in-<k>   (process input,  LEFT  of pair)
% Time arrow:  <name>-t0  (right)  -->  <name>-tn (left)
\newcommand{\PTDrawProcessMPO}[1][]{%
  \PTMPOjointfalse
  \PTMPOtagstrue
  \PTMPOchitrue
  \PTMPOtimestrue
  \PTMPOarrowtrue
  \PTMPOfinalfalse
  \setlength{\PTMPOspacing}{\PTTimeSpacing}%
  \def\PTMPOn{3}%
  \def\PTMPOat{0,0}%
  \def\PTMPOname{pt}%
  \def\PTMPOlabel{}%
  \setkeys{ptmpo}{#1}%
  \coordinate (\PTMPOname-origin) at (\PTMPOat);
  \edef\PTMPOlast{\PTMPOname-W-\the\numexpr\PTMPOn-1\relax}%
  % --- cores, right (k=0) to left ------------------------------------------
  \foreach \k in {0,...,\the\numexpr\PTMPOn-1\relax} {%
    \path (\PTMPOname-origin)
      ++(-\k*\PTMPOspacing,0)
      coordinate (\PTMPOname-c-\k);
    \PTProcessCore{\PTMPOname-W-\k}{(\PTMPOname-c-\k)}{\(W^{[\k]}\)}%
    % Local ports: output on the right, input on the left (local time).
    \coordinate (\PTMPOname-outsrc-\k) at
      ([xshift=\PTIOSep]\PTMPOname-W-\k.south);
    \coordinate (\PTMPOname-insrc-\k) at
      ([xshift=-\PTIOSep]\PTMPOname-W-\k.south);
    \draw[PTLiouvilleLeg]
      (\PTMPOname-outsrc-\k) -- ++(0,-\PTLegLength)
      coordinate (\PTMPOname-out-\k);
    \draw[PTLiouvilleLeg]
      (\PTMPOname-insrc-\k) -- ++(0,-\PTLegLength)
      coordinate (\PTMPOname-in-\k);
  }%
  % --- joint body (no exposed link indices) or MPO links -------------------
  \ifPTMPOjoint
    \node[
      fit=(\PTMPOname-W-0) (\PTMPOlast),
      inner xsep=3.6pt,
      inner ysep=4.6pt,
      rounded corners=5.5pt,
      fill=PTProcess!55,
      draw=black,
      line width=0.7pt,
    ] (\PTMPOname-body) {};
    \draw[
      draw=PTProcess!70!black,
      line width=3.4pt,
      line cap=round,
    ] (\PTMPOlast.center) -- (\PTMPOname-W-0.center);
  \else
    \ifnum\PTMPOn>1
      \foreach \k [
        evaluate=\k as \knext using int(\k+1)
      ] in {0,...,\the\numexpr\PTMPOn-2\relax} {%
        \draw[PTLinkLeg]
          (\PTMPOname-W-\k.west) -- (\PTMPOname-W-\knext.east)
          coordinate[midway] (\PTMPOname-link-\k);
        \ifPTMPOchi
          \node[font=\small, text=PTLabel, inner sep=0.5pt, above=1.4pt]
            at (\PTMPOname-link-\k) {\(\chi_{\k}\)};
        \fi
        \ifPTMPOtags
          \PTIndexLabel[
            anchor=south,
            tag={PTLink,t-0\k},
            yshift=11.0pt,
          ]{\PTMPOname-link-\k}%
        \fi
      }%
    \fi
    % Dimension-1 boundary stubs.
    \draw[PTLinkLeg] (\PTMPOname-W-0.east) -- ++(4.2mm,0)
      coordinate (\PTMPOname-bnd-R);
    \begin{scope}[on PTforeground layer]
      \node[PTDummyBond] (\PTMPOname-dummy-R) at (\PTMPOname-bnd-R) {};
    \end{scope}
    \draw[PTLinkLeg]
      (\PTMPOlast.west) -- ++(-4.2mm,0)
      coordinate (\PTMPOname-bnd-L);
    \begin{scope}[on PTforeground layer]
      \node[PTDummyBond] (\PTMPOname-dummy-L) at (\PTMPOname-bnd-L) {};
    \end{scope}
  \fi
  % --- optional final (t_n) physical output on the far left ----------------
  \ifPTMPOfinal
    \coordinate (\PTMPOname-final-src) at ([yshift=1.6mm]\PTMPOlast.west);
    \draw[PTLiouvilleLeg] (\PTMPOname-final-src) -- ++(-7.5mm,0)
      coordinate (\PTMPOname-final);
  \fi
  % --- time labels and left-pointing arrow ---------------------------------
  \coordinate (\PTMPOname-t0) at ([yshift=-15.2mm, xshift=5.0mm]\PTMPOname-W-0.south);
  \coordinate (\PTMPOname-tn) at ([yshift=-15.2mm, xshift=-5.0mm]\PTMPOlast.south);
  \ifPTMPOtimes
    \node[font=\small, text=PTLabel, inner sep=0.4pt, anchor=west]
      at (\PTMPOname-t0) {\(t_{0}\)};
    \node[font=\small, text=PTLabel, inner sep=0.4pt, anchor=east]
      at (\PTMPOname-tn) {\(t_{n}\)};
  \fi
  \ifPTMPOarrow
    \PTTimeArrow[below=1.2pt]{\PTMPOname-t0}{\PTMPOname-tn}{\(t\)}%
  \fi
  \ifPTMPOtags
    \PTIndexLabel[
      anchor=west,
      math={\ell_{\mathrm{out}}},
      tag={Output,t-00},
      xshift=3.0pt,
    ]{\PTMPOname-out-0}%
    \PTIndexLabel[
      anchor=east,
      math={\ell_{\mathrm{in}}},
      tag={Input,t-00},
      xshift=-3.0pt,
      yshift=-4.2pt,
    ]{\PTMPOname-in-0}%
  \fi
  \def\PTMPOlabelEmpty{}%
  \ifx\PTMPOlabel\PTMPOlabelEmpty
  \else
    \node[font=\small, text=PTLabel, inner sep=1pt, above=3.6pt]
      at ($ (\PTMPOname-W-0.north)!0.5!(\PTMPOlast.north) $)
      {\PTMPOlabel};
  \fi
}

\newcommand{\PTDrawJointProcess}[1][]{%
  \PTDrawProcessMPO[joint=true,chi=false,tags=false,#1]%
}

% -----------------------------------------------------------------------------
% Example of a future instrument (commented; copy and adapt).
% -----------------------------------------------------------------------------
% \newcommand{\PTReset}[4][]{%
%   \PTInstrumentNode[fill=PTInstrument, classical=none, #1]{#2}{#3}{#4}%
% }
% % Tester / ancilla interaction (two-site box):
% \newcommand{\PTTester}[4][]{%
%   \PTInstrumentNode[
%     fill=PTInstrument,
%     minimum width=1.55\PTCoreWidth,
%     #1
%   ]{#2}{#3}{#4}%
% }
% % Feedback-conditioned operation: reuse classical=down and wire it into
% % a later \PTInstrumentNode.
%
% This file is a library.  Compile pt_tikz_demo.tex for a catalogue of
% every primitive.
